\documentclass[11pt,a4paper,oneside]{article}

\usepackage[T1]{fontenc}
\usepackage[utf8]{inputenc}
\usepackage[brazil]{babel}
\usepackage{lmodern}
\usepackage{microtype}
\usepackage{geometry}
\usepackage{setspace}
\usepackage{enumitem}
\usepackage{booktabs}
\usepackage{longtable}
\usepackage{array}
\usepackage{tabularx}
\usepackage{ltablex}
\usepackage{xcolor}
\usepackage{hyperref}
\usepackage{fancyhdr}
\usepackage{titlesec}
\usepackage{tcolorbox}
\usepackage{tikz}
\usepackage{pdflscape}
\usepackage{listings}
\usepackage{caption}
\usepackage{ragged2e}
\usetikzlibrary{arrows.meta,positioning,fit,shapes.geometric}

\geometry{top=2.2cm,bottom=2.2cm,left=2.2cm,right=2.2cm}
\setstretch{1.08}
\setlength{\parindent}{0pt}
\setlength{\parskip}{0.55em}
\setlist[itemize]{leftmargin=1.45em,itemsep=0.15em,topsep=0.25em}
\setlist[enumerate]{leftmargin=1.65em,itemsep=0.2em,topsep=0.25em}
\renewcommand{\arraystretch}{1.18}
\keepXColumns

\definecolor{sisterblue}{RGB}{28,67,105}
\definecolor{sisterlightblue}{RGB}{232,241,248}
\definecolor{sistergray}{RGB}{245,246,247}
\definecolor{sisterdarkgray}{RGB}{74,83,91}
\definecolor{sistergreen}{RGB}{33,112,77}
\definecolor{sisteramber}{RGB}{161,102,0}
\definecolor{sisterred}{RGB}{155,44,44}
\definecolor{sisterpurple}{RGB}{93,63,130}

\hypersetup{
  colorlinks=true,
  linkcolor=sisterblue,
  urlcolor=sisterblue,
  citecolor=sisterblue,
  pdftitle={Especificacao Funcional e de Engenharia do SisTer},
  pdfauthor={Projeto SisTer},
  pdfsubject={Arquitetura funcional, rastreabilidade e materializacao incremental em codigo}
}

\pagestyle{fancy}
\fancyhf{}
\fancyhead[L]{\small\textcolor{sisterdarkgray}{SisTer -- Especificação Funcional e de Engenharia}}
\fancyhead[R]{\small\textcolor{sisterdarkgray}{EFE-SisTer/1.4}}
\fancyfoot[C]{\small Página \thepage}
\renewcommand{\headrulewidth}{0.3pt}
\renewcommand{\footrulewidth}{0pt}

\titleformat{\section}{\Large\bfseries\color{sisterblue}}{\thesection.}{0.55em}{}
\titleformat{\subsection}{\large\bfseries\color{sisterdarkgray}}{\thesubsection.}{0.5em}{}
\titleformat{\subsubsection}{\normalsize\bfseries}{\thesubsubsection.}{0.45em}{}

\newcolumntype{Y}{>{\RaggedRight\arraybackslash}X}
\newcolumntype{P}[1]{>{\RaggedRight\arraybackslash}p{#1}}

\newtcolorbox{principio}[1][]{
  colback=sisterlightblue,
  colframe=sisterblue,
  boxrule=0.65pt,
  arc=1.5mm,
  left=3.5mm,right=3.5mm,top=2.2mm,bottom=2.2mm,
  title={#1},fonttitle=\bfseries
}
\newtcolorbox{decisao}[1][]{
  colback=sistergray,
  colframe=sisterdarkgray,
  boxrule=0.55pt,
  arc=1.5mm,
  left=3.5mm,right=3.5mm,top=2mm,bottom=2mm,
  title={#1},fonttitle=\bfseries
}
\newtcolorbox{alerta}[1][]{
  colback=red!3,
  colframe=sisterred,
  boxrule=0.6pt,
  arc=1.5mm,
  left=3.5mm,right=3.5mm,top=2mm,bottom=2mm,
  title={#1},fonttitle=\bfseries
}

\newcommand{\statusimpl}{\textcolor{sistergreen}{\textbf{IMPLEMENTADO}}}
\newcommand{\statusparcial}{\textcolor{sisteramber}{\textbf{PARCIAL}}}
\newcommand{\statusfallback}{\textcolor{sisterpurple}{\textbf{FALLBACK}}}
\newcommand{\statusespec}{\textcolor{sisterblue}{\textbf{ESPECIFICADO}}}
\newcommand{\statusausente}{\textcolor{sisterred}{\textbf{NÃO OBSERVADO}}}
\newcommand{\statusnv}{\textcolor{sisterdarkgray}{\textbf{NÃO VERIFICÁVEL}}}

\lstdefinestyle{sistercode}{
  basicstyle=\ttfamily\small,
  backgroundcolor=\color{sistergray},
  frame=single,
  rulecolor=\color{sisterdarkgray!35},
  breaklines=true,
  columns=fullflexible,
  keepspaces=true,
  showstringspaces=false,
  xleftmargin=0.5em,
  framexleftmargin=0.4em
}
\lstset{style=sistercode}

\begin{document}

\begin{titlepage}
\pagestyle{empty}
\thispagestyle{empty}
\vspace*{2.2cm}
{\color{sisterblue}\rule{\textwidth}{1.2pt}}\\[1.2cm]
{\Huge\bfseries Especificação Funcional e de Engenharia do SisTer\par}
\vspace{0.45cm}
{\Large Do objetivo científico à implementação, reflexividade, evidência e maturidade\par}
\vspace{1.2cm}
\begin{principio}[Pergunta governante]
\centering\Large\textbf{Como o SisTer realiza seu propósito, avalia a própria execução e se reconfigura de modo governado desde a primeira decisão de engenharia até a produção de conhecimento rastreável?}
\end{principio}
\vfill
\begin{tabularx}{\textwidth}{P{4.2cm}Y}
\toprule
\textbf{Identificador} & EFE-SisTer/1.4 \\
\textbf{Status} & Versão revisada para validação; formaliza níveis, perfis, autoridade e materialização C++ da reflexividade operacional \\
\textbf{Linha de base do código} & Snapshot \texttt{apps.zip}, recebido em 1º de agosto de 2026 \\
\textbf{Escopo verificado} & \texttt{apps/sisterd} e \texttt{apps/sisterctl} \\
\textbf{Processo de referência} & Relação Nexo--Compras: dado contratado $\rightarrow$ informação integrada $\rightarrow$ conhecimento tipificado \\
\textbf{Data} & 2 de agosto de 2026 \\
\bottomrule
\end{tabularx}
\vspace{1.2cm}
{\color{sisterblue}\rule{\textwidth}{1.2pt}}
\thispagestyle{empty}
\end{titlepage}
\pagestyle{fancy}

\pagenumbering{roman}
\section*{Controle do documento}
\addcontentsline{toc}{section}{Controle do documento}

\begin{tabularx}{\textwidth}{P{3.6cm}Y}
\toprule
\textbf{Campo} & \textbf{Definição} \\
\midrule
Natureza & Especificação técnica de engenharia: normativa no modelo-alvo e descritiva na linha de base verificada. \\
Autoridade & Documento-mestre de arquitetura funcional, subordinado à finalidade e à ontologia do SisTer e articulador dos planos, contratos, ADRs e perfis de maturidade. \\
Não substitui & ADRs de decisões específicas; contratos e schemas; manuais operacionais; planos de implantação; evidências produzidas pelo SGE-SisTer. \\
Deve governar & A relação entre objetivo, função, processo, capacidade, contrato, componente, código, teste, evidência, avaliação reflexiva, ação corretiva e maturidade. \\
Regra de atualização & Toda mudança que altere função sistêmica, processo finalístico, fronteira, contrato estrutural ou cadeia de rastreabilidade deve atualizar esta especificação ou registrar explicitamente por que não a altera. \\
\bottomrule
\end{tabularx}

\subsection*{Histórico de revisão}
\begin{tabularx}{\textwidth}{P{2.0cm}P{2.6cm}Y}
\toprule
\textbf{Versão} & \textbf{Data} & \textbf{Alteração} \\
\midrule
1.0 & 01/08/2026 & Especificação inicial. \\
1.1 & 01/08/2026 & Correção de estados, dependências, ciclos de vida, rastreabilidade, processo Nexo--Compras e roadmap; redução de conteúdo não normativo. \\
1.2 & 01/08/2026 & Incorporação da engenharia de segurança orientada por ameaças, perfil restrito de transporte, classes iniciais de ameaça, rastreabilidade de controles e ciclo contínuo de resposta a vulnerabilidades. \\
1.3 & 02/08/2026 & Incorporação do Princípio da Reflexividade Operacional, da função F-11, do processo P-08 e dos artefatos de avaliação e ação corretiva governada. \\
1.4 & 02/08/2026 & Formalização das profundidades D0--D5, autoridades A0--A5, efeitos de gate, respostas operacionais, \texttt{ReflexivityProfile}, arquitetura C++ e piloto Nexo--Compras em \texttt{shadow}. \\
\bottomrule
\end{tabularx}

\subsection*{Fontes consolidadas}
Esta versão foi construída a partir das fontes e linhas de elaboração disponibilizadas para análise:
\begin{itemize}
  \item \texttt{sister\_funcoes\_processos\_atividade\_cientifica(2).tex};
  \item \texttt{plano\_transicao\_sister\_engenharia\_governada\_revisado(2).md};
  \item \texttt{plano\_trabalho\_sister\_kanban(2).md};
  \item \texttt{funcoes\_sistemicas\_sisterd\_sisterctl(1).md};
  \item snapshot de código \texttt{apps.zip};
  \item achados reproduzidos de robustez do transporte e avaliação das decisões SEC-00 e SEC-01;
  \item reflexão arquitetural sobre cibernética, retroalimentação e avaliação da própria execução do SisTer;
  \item avaliação dos níveis de profundidade, autoridade e efeito operacional da reflexividade, incluindo diretrizes de materialização com C++20/23 e laboratório C++26;
  \item referências normativas e históricas de segurança listadas no Apêndice~\ref{sec:referencias-seguranca}.
\end{itemize}

\begin{alerta}[Limite de verificação]
A linha de base de implementação não constitui auditoria do repositório completo. Foram inspecionados 3.737 linhas distribuídas em \texttt{apps/sisterd} e \texttt{apps/sisterctl}. Contratos, migrations, scripts, biblioteca \texttt{sister/contract.hpp}, interfaces web e testes externos ao arquivo \texttt{apps.zip} são mencionados apenas quando sustentados pelos documentos fornecidos; sua implementação não foi verificada neste snapshot.
\end{alerta}

\tableofcontents
\clearpage
\pagenumbering{arabic}

\section{Finalidade e tese de engenharia}

\subsection{Finalidade científica}
\begin{principio}[Definição central]
\textbf{O SisTer é uma plataforma federativa que articula subsistemas autônomos, governa seus relacionamentos por contratos e transforma ofertas de diferentes domínios em informação integrada e conhecimento contextualizado, preservando autonomia, proveniência, responsabilidade e capacidade de revisão.}
\end{principio}

O sistema existe para apoiar a caracterização, o monitoramento e a compreensão da resiliência de Sistemas Socioecológicos (SSE), inclusive em Sistemas Agroalimentares de Interesse Territorial (SAIT), sem substituir a autoridade científica e operacional dos domínios participantes.

\subsection{Tese de engenharia}
\begin{principio}[Regra de materialização]
Toda implementação governada deve ser justificável como materialização de uma função necessária ao objetivo do SisTer, executada em um processo definido, limitada por contrato e comprovada por testes e evidências.
\end{principio}

A unidade de evolução é a \textbf{capacidade rastreável}, com finalidade, responsável, condições de uso, critérios de aceitação e evidência.

\subsection{Princípio da Reflexividade Operacional}
\label{sec:reflexividade-operacional}
\begin{principio}[Princípio da Reflexividade Operacional]
\textbf{Toda função relevante do SisTer deve produzir evidências suficientes para que o sistema possa avaliar sua própria execução em relação ao propósito, aos contratos e às políticas vigentes, tornando divergências compreensíveis e ações corretivas governáveis.}
\end{principio}

Reflexividade operacional não significa consciência, intenção própria ou liberdade para alterar autonomamente o código, os contratos ou os critérios científicos. Significa dotar cada execução relevante de referências explícitas, evidências observáveis e um mecanismo versionado de comparação capaz de responder:
\begin{itemize}
  \item o que deveria ter ocorrido e sob qual propósito, contrato, política e revisão;
  \item o que efetivamente ocorreu, com quais entradas, decisões, transformações e resultados;
  \item em que medida a execução confirmou, divergiu ou permaneceu inconclusiva diante das referências;
  \item qual resposta é admissível, quem possui autoridade para autorizá-la e como seu efeito será novamente avaliado.
\end{itemize}

\begin{lstlisting}
proposito + contrato + politica + criterio
-> execucao
-> evidencias
-> avaliacao reflexiva
-> confirmacao | divergencia | inconclusao
-> decisao autorizada
-> correcao | reconfiguracao | revisao da referencia
-> nova execucao e nova avaliacao
\end{lstlisting}

\begin{decisao}[Limite de autonomia]
O SisTer possui autonomia para observar, comparar, explicar e recomendar. Somente pode executar ações corretivas quando elas estiverem previamente autorizadas por política, forem proporcionais ao risco, auditáveis e reversíveis. Mudanças estruturais em contratos, critérios científicos, políticas de maturidade ou código exigem decisão humana ou processo formal equivalente.
\end{decisao}

A confirmação também é resultado relevante: ela sustenta confiança sem transformar ausência de alerta em prova de correção. A inconclusão deve permanecer explícita quando faltarem evidências, referências ou capacidade avaliativa.


\subsection{Modelo normativo de aplicação da reflexividade}
\label{sec:modelo-reflexividade}

A reflexividade não possui um único nível global. Cada capacidade relevante deve declarar um \texttt{ReflexivityProfile} que combine três dimensões independentes:

\begin{center}
\begin{tikzpicture}[
  node distance=8mm,
  box/.style={draw=sisterpurple, rounded corners, fill=sisterpurple!8, minimum width=4.4cm, minimum height=12mm, align=center, font=\small\bfseries},
  arr/.style={-{Latex[length=2.2mm]}, thick, draw=sisterdarkgray}
]
\node[box] (d) {Profundidade\\o que é avaliado};
\node[box, right=of d] (a) {Autoridade\\o que o sistema pode fazer};
\node[box, right=of a] (g) {Efeito operacional\\como a conclusão afeta o fluxo};
\draw[arr] (d)--(a); \draw[arr] (a)--(g);
\end{tikzpicture}
\end{center}

\begin{decisao}[Regra de composição]
Profundidade, autoridade e efeito operacional não podem ser inferidos um do outro. Uma avaliação profunda pode permanecer apenas consultiva; uma proteção automática pode ser legítima para uma invariante estrutural simples; e uma divergência não implica necessariamente bloqueio ou correção.
\end{decisao}

\subsubsection{Reflexividade do SisTer e reflexão da linguagem}
A \textbf{reflexividade operacional do SisTer} avalia uma execução em relação ao propósito, aos contratos, às políticas e aos critérios aplicáveis. A \textbf{reflexão estática do C++} inspeciona estruturas do programa em tempo de compilação e pode apoiar geração de metadados, schemas, serialização e validação. São conceitos complementares, mas não equivalentes:

\begin{principio}[Distinção obrigatória]
\textbf{O C++ pode fornecer mecanismos para o sistema observar sua estrutura; somente o modelo funcional, as referências e as evidências permitem ao SisTer avaliar o significado de sua execução.}
\end{principio}

\subsubsection{Profundidades D0--D5}
\begin{longtable}{P{1.5cm}P{3.4cm}P{5.0cm}P{5.0cm}}
\toprule
\textbf{Nível} & \textbf{Nome} & \textbf{Objeto da avaliação} & \textbf{Aplicação no SisTer} \\
\midrule
\endhead
D0 & Estrutural & Tipos, interfaces, estados e invariantes verificáveis em compilação ou validação local. & Impedir mistura de identificadores, estados impossíveis e contratos internos incompletos. \\
D1 & Funcional & Resultado de uma função, comando ou operação isolada. & Verificar precondições, pós-condições, erro tipificado, evidência e cancelamento de uma operação. \\
D2 & Processual & Ciclo completo de uma execução governada. & Avaliar se uma \texttt{IntegrationRun} percorreu estados válidos e produziu resultados e evidências suficientes. \\
D3 & Relacional & Interação entre participantes autônomos. & Avaliar se capacidades, versões, finalidade, restrições e responsabilidades do acordo foram respeitadas. \\
D4 & Ecossistêmica & Coerência entre componentes, políticas, dependências e maturidade. & Detectar regressões, conflitos de política, degradação sistêmica e impacto entre capacidades. \\
D5 & Epistêmica & Adequação dos critérios, métodos ou referências usados para avaliar. & Sinalizar contradições, inconclusões recorrentes ou necessidade de revisão científica/institucional, sem alterar a referência autonomamente. \\
\bottomrule
\end{longtable}

Aplicação mínima adotada:
\begin{itemize}
  \item D0--D1 em toda capacidade governada;
  \item D2 em toda execução relevante;
  \item D3 em todo relacionamento contratado;
  \item D4 no SGE-SisTer e no plano de controle;
  \item D5 somente com owner científico ou institucional explícito e sem revisão automática da referência.
\end{itemize}

\subsubsection{Autoridades A0--A5}
\begin{longtable}{P{1.5cm}P{3.7cm}P{7.0cm}P{2.8cm}}
\toprule
\textbf{Nível} & \textbf{Autoridade} & \textbf{Limite operacional} & \textbf{Modo típico} \\
\midrule
\endhead
A0 & Observar & Produzir eventos, métricas, estados e evidências, sem emitir juízo reflexivo. & \texttt{observe} \\
A1 & Avaliar & Comparar evidências e referências, produzindo resultado explicável sem alterar o fluxo. & \texttt{shadow} \\
A2 & Recomendar & Explicar causa, impacto, alternativas, owner e ação admissível, sem autorização implícita. & \texttt{advisory} \\
A3 & Proteger ou bloquear & Impedir operação incompatível antes de efeitos indevidos, conforme política explícita. & \texttt{enforcing} \\
A4 & Corrigir dentro de limites & Executar somente ação previamente autorizada, proporcional, auditável, reversível e reavaliada. & \texttt{bounded\_correction} \\
A5 & Revisar referência & Alterar propósito, contrato, política, critério ou método por decisão humana ou processo formal equivalente. & governança humana \\
\bottomrule
\end{longtable}

A5 não constitui autonomia de software. O SisTer pode produzir uma \texttt{ReferenceRevisionProposal}, mas a aprovação pertence ao owner funcional, científico, arquitetural ou institucional da referência.

\subsubsection{Resultado, gate e resposta operacional}
Três vocabulários devem permanecer separados:

\begin{tabularx}{\textwidth}{P{4.0cm}Y}
\toprule
\textbf{Dimensão} & \textbf{Valores mínimos} \\
\midrule
Resultado da avaliação & \texttt{confirmed}, \texttt{divergent}, \texttt{inconclusive}, \texttt{not\_applicable}. \\
Efeito sobre gate/operação & \texttt{shadow}, \texttt{pass}, \texttt{warn}, \texttt{block}. \\
Resposta operacional & \texttt{none}, \texttt{collect\_evidence}, \texttt{recommend}, \texttt{request\_approval}, \texttt{cancel}, \texttt{retry}, \texttt{isolate}, \texttt{rollback}, \texttt{suspend}, \texttt{reprocess}. \\
\bottomrule
\end{tabularx}

Uma avaliação \texttt{divergent} pode gerar apenas \texttt{warn} e recomendação. Outra, diante de invariante crítica e política A3, pode gerar \texttt{block}. O resultado descreve a relação com a referência; o gate determina o efeito; a ação responde ao efeito autorizado.

\subsubsection{Perfil inicial por escopo}
\begin{longtable}{P{5.0cm}P{2.3cm}P{3.0cm}P{4.4cm}}
\toprule
\textbf{Escopo} & \textbf{Profundidade} & \textbf{Autoridade inicial} & \textbf{Destino permitido} \\
\midrule
\endhead
Invariantes internas do C++ & D0--D1 & A3 & Permanecer em A3; falha cedo e explícita. \\
Execução Nexo--Compras & D2 & A1 & Evoluir para A2/A3 após histórico e critérios aceitos. \\
Acordo Nexo--Compras & D3 & A1--A2 & Evoluir para A3 em invariantes contratuais objetivas. \\
Saúde de adaptadores & D1--D2 & A2 & A4 limitado a cancelamento, isolamento e recuperação autorizados. \\
Maturidade pelo SGE-SisTer & D4 & A1--A2 & A3 para promoção, regressão e publicação governadas. \\
Correlação e interpretação científica & D3--D5 & A1--A2 & Decisão humana para revisão de método ou significado. \\
Alteração de contrato, política, método ou código & -- & nenhuma & A5 por processo formal de governança e engenharia. \\
\bottomrule
\end{longtable}

\begin{principio}[Diretriz adotada]
\centering
\textbf{O SisTer deve nascer amplamente reflexivo, porém apenas seletivamente autocorretivo.}
\end{principio}

\subsubsection{Decisão para o primeiro piloto}
O processo Nexo--Compras deve iniciar com profundidade D2--D3, autoridade A1 e efeito \texttt{shadow}. O sistema congela referências, captura evidências, produz \texttt{OperationalAssessment} e explica a conclusão pelo \texttt{sisterctl}, sem bloquear ou corrigir a execução. A passagem a A2, A3 ou A4 depende de evidência histórica, ADR aceita, política explícita, testes de falha e comprovação de reversibilidade.

\subsection{Princípios estruturantes}
\begin{enumerate}
  \item \textbf{Autonomia federativa}: cada subsistema preserva domínio, dados de autoridade, regras, banco e responsabilidade.
  \item \textbf{Contratos nas fronteiras}: integração não decorre de conhecimento informal do código interno.
  \item \textbf{Convergência sem centralização indevida}: o SisTer relaciona contribuições, mas não absorve os domínios.
  \item \textbf{Transformação epistemicamente rastreável}: dado, informação e conhecimento possuem critérios distintos e relações reproduzíveis.
  \item \textbf{Experimentação livre na borda}: hipóteses podem nascer em laboratórios e protótipos.
  \item \textbf{Evidência obrigatória na promoção}: somente capacidades comprovadas podem integrar, governar ou publicar.
  \item \textbf{Governança proporcional ao risco}: os controles crescem conforme impacto, criticidade e exposição.
  \item \textbf{Segurança orientada por ameaças}: capacidades, fronteiras e protocolos devem herdar o conhecimento acumulado sobre riscos aplicáveis e demonstrar seus controles por evidências.
  \item \textbf{Reflexividade operacional}: toda execução relevante produz evidências e é comparável ao propósito, aos contratos, às políticas e aos critérios vigentes.
  \item \textbf{Reversibilidade}: maturidade, contrato e conhecimento podem ser suspensos, invalidados, revogados ou reprocessados.
\end{enumerate}

\section{Vocabulário normativo e hierarquia}

\subsection{Termos fundamentais}
\begin{longtable}{P{3.2cm}P{5.6cm}P{6.0cm}}
\toprule
\textbf{Termo} & \textbf{Pergunta respondida} & \textbf{Definição operacional} \\
\midrule
\endhead
Objetivo & Por que o SisTer existe? & Resultado científico, institucional e social que justifica o sistema. \\
Função & O que o sistema precisa fazer continuamente? & Responsabilidade sistêmica estável, independente de um componente específico. \\
Processo & Como funções, atores e artefatos se coordenam ao longo do tempo? & Fluxo governado com gatilho, estados, participantes, entradas, saídas, falhas e evidências. \\
Capacidade & Há meios para realizar determinada função? & Unidade técnica, organizacional e contratual que pode ser oferecida, autorizada, testada e avaliada. \\
Serviço & Por qual interface a capacidade é disponibilizada? & API, evento, pacote, comando, consulta, adaptador ou procedimento humano governado. \\
Componente & Onde a capacidade está implementada? & Unidade de software ou infraestrutura sob responsabilidade definida. \\
Contrato & Sob quais condições participantes podem interagir? & Identidade, semântica, versão, finalidade, autoridade, restrições, estados e evidências da relação. \\
Evidência & Como sabemos que uma afirmação operacional é verdadeira? & Registro estruturado, versionado, associável a execução, revisão, ambiente, resultado e artefatos. \\
Referência de avaliação & Em relação a que a execução será julgada? & Snapshot versionado do propósito, contrato, política, critério de aceitação, risco e contexto aplicáveis à execução. \\
Avaliação reflexiva & A execução permaneceu coerente com suas referências? & Comparação reproduzível entre evidências observadas e referências vigentes, com resultado, divergências, limitações, responsável e recomendação. \\
Ação corretiva & O que pode ser feito diante da divergência? & Resposta autorizada, proporcional, auditável e preferencialmente reversível, com escopo, responsável, condição de aplicação e verificação posterior. \\
Perfil de reflexividade & Qual avaliação e autoridade se aplicam a esta capacidade? & Contrato operacional que declara escopo, profundidade, modo, referências, evidências, avaliador, tolerâncias, efeitos de gate, ações permitidas, owner, reversão, retenção e política de falha. \\
Profundidade reflexiva & Até que camada a avaliação alcança? & Nível D0--D5 que distingue avaliação estrutural, funcional, processual, relacional, ecossistêmica e epistêmica. \\
Autoridade reflexiva & O que o sistema está autorizado a fazer? & Nível A0--A5 que distingue observar, avaliar, recomendar, proteger, corrigir dentro de limites e revisar referência por governança humana. \\
Efeito de gate & Como a conclusão afeta operação ou maturidade? & Resultado independente da avaliação, expresso por \texttt{shadow}, \texttt{pass}, \texttt{warn} ou \texttt{block}. \\
Ameaça & O que pode comprometer um ativo ou garantia? & Cenário identificável, conhecido ou plausível, associado a superfície, precondições, impacto e exposição. \\
Controle & Como a ameaça é prevenida, detectada, contida ou recuperada? & Medida técnica, operacional ou de governança com owner, teste e evidência. \\
Risco residual & O que permanece após os controles? & Exposição remanescente explicitamente aceita, bloqueada ou limitada por maturidade e contexto operacional. \\
Maturidade & Qual confiança é autorizada para esta capacidade? & Estado governado, reversível, sustentado por evidências e limitado por perfil e risco. \\
\bottomrule
\end{longtable}

\subsection{Hierarquia obrigatória}
\begin{center}
\begin{tikzpicture}[
  node distance=5mm,
  box/.style={draw=sisterblue, rounded corners, fill=sisterlightblue, minimum width=8.4cm, minimum height=8mm, align=center, font=\small\bfseries},
  arr/.style={-{Latex[length=2.2mm]}, thick, draw=sisterdarkgray}
]
\node[box] (o) {Objetivo e ontologia};
\node[box, below=of o] (f) {Funções sistêmicas};
\node[box, below=of f] (p) {Processos finalísticos};
\node[box, below=of p] (c) {Capacidades};
\node[box, below=of c] (ct) {Contratos e acordos};
\node[box, below=of ct] (s) {Serviços e componentes};
\node[box, below=of s] (i) {Implementação};
\node[box, below=of i] (e) {Testes, evidências, avaliação reflexiva e maturidade};
\draw[arr] (o)--(f); \draw[arr] (f)--(p); \draw[arr] (p)--(c); \draw[arr] (c)--(ct);
\draw[arr] (ct)--(s); \draw[arr] (s)--(i); \draw[arr] (i)--(e);
\draw[arr, bend right=68] (e.west) to node[left, font=\scriptsize, align=right]{retroalimentação\\governada} (f.west);
\end{tikzpicture}
\end{center}

\begin{decisao}[Regra contra inversão arquitetural]
Classes, endpoints, scripts e tabelas não definem a finalidade do sistema. Eles somente são admitidos como elementos governados quando podem ser remontados à hierarquia acima.
\end{decisao}

\section{Processo finalístico completo}

\subsection{Ciclo de transformação e retroalimentação}
\begin{center}
\resizebox{\textwidth}{!}{%
\begin{tikzpicture}[
  node distance=7mm,
  stage/.style={draw=sisterblue, rounded corners, fill=sisterlightblue, minimum width=2.25cm, minimum height=10mm, align=center, font=\small\bfseries},
  arr/.style={-{Latex[length=2.3mm]}, thick, draw=sisterdarkgray}
]
\node[stage] (r) {Realidade};
\node[stage, right=of r] (o) {Observação};
\node[stage, right=of o] (d) {Dado};
\node[stage, right=of d] (v) {Validação};
\node[stage, right=of v] (i) {Integração};
\node[stage, right=of i] (inf) {Informação};
\node[stage, right=of inf] (c) {Convergência};
\node[stage, right=of c] (k) {Conhecimento};
\node[stage, right=of k] (dec) {Decisão};
\node[stage, right=of dec] (a) {Ação};
\draw[arr] (r)--(o); \draw[arr] (o)--(d); \draw[arr] (d)--(v); \draw[arr] (v)--(i);
\draw[arr] (i)--(inf); \draw[arr] (inf)--(c); \draw[arr] (c)--(k); \draw[arr] (k)--(dec); \draw[arr] (dec)--(a);
\draw[arr, bend left=35] (a.north) to node[above, font=\scriptsize]{nova observação} (r.north);
\end{tikzpicture}}
\end{center}

O ciclo completo é maior do que o software central. A observação pode ocorrer em MorfoCampo, Clima, Compras, Nexo ou outro domínio. O SisTer governa especialmente as passagens entre fronteiras: identidade, autorização, contrato, validação, integração, proveniência, convergência, exposição e revisão.

\subsection{Ciclo reflexivo da execução governada}
O ciclo epistemológico descreve como a ação retorna à realidade observada. O ciclo reflexivo descreve como o próprio SisTer observa o desempenho de suas funções e verifica se a execução permaneceu coerente com aquilo que a autorizou.

\begin{center}
\resizebox{0.98\textwidth}{!}{%
\begin{tikzpicture}[
  node distance=7mm,
  stage/.style={draw=sisterpurple, rounded corners, fill=sisterpurple!8, minimum width=2.7cm, minimum height=10mm, align=center, font=\small\bfseries},
  arr/.style={-{Latex[length=2.3mm]}, thick, draw=sisterdarkgray}
]
\node[stage] (ref) {Referências\\vigentes};
\node[stage, right=of ref] (exe) {Execução};
\node[stage, right=of exe] (evi) {Evidências};
\node[stage, right=of evi] (ava) {Avaliação\\reflexiva};
\node[stage, right=of ava] (res) {Confirmação,\\divergência ou\\inconclusão};
\node[stage, right=of res] (dec) {Decisão\\autorizada};
\node[stage, right=of dec] (cor) {Correção ou\\reconfiguração};
\draw[arr] (ref)--(exe); \draw[arr] (exe)--(evi); \draw[arr] (evi)--(ava);
\draw[arr] (ava)--(res); \draw[arr] (res)--(dec); \draw[arr] (dec)--(cor);
\draw[arr, bend left=34] (cor.north) to node[above, font=\scriptsize]{nova execução} (exe.north);
\draw[arr, dashed, bend right=38] (res.south) to node[below, font=\scriptsize]{descoberta pode revisar a referência por processo formal} (ref.south);
\end{tikzpicture}}
\end{center}

A avaliação deve registrar a revisão exata das referências utilizadas. Caso contrário, uma política alterada posteriormente poderia reescrever retrospectivamente o significado de uma execução passada. O resultado mínimo é \texttt{confirmed}, \texttt{divergent}, \texttt{inconclusive} ou \texttt{not\_applicable}; severidade e efeito sobre operação e maturidade são dimensões separadas.

\subsection{Critérios das passagens epistemológicas}
\begin{longtable}{P{3.0cm}P{5.4cm}P{6.4cm}}
\toprule
\textbf{Estágio} & \textbf{Condição mínima} & \textbf{Evidência exigida} \\
\midrule
\endhead
Dado contratado & Produtor identificado, contrato ativo, schema e versão reconhecidos, finalidade autorizada, integridade verificável. & Identificador imutável, digest, instante, produtor, contrato, schema, escopo e referência ao original. \\
Informação integrada & Dado contextualizado, associado a entidades, tempo e espaço, relacionado a outras fontes, com incertezas explícitas. & Transformação identificada, entradas, parâmetros, código/revisão, relações produzidas, pendências e qualidade. \\
Conhecimento rastreável & Afirmação ou interpretação sustentada por informações, método, responsabilidade, limitações e confiança. & \texttt{KnowledgeArtifact} revisável, autoria, evidências, método, validade, confiança, finalidade e histórico. \\
Decisão governada & Conhecimento utilizado por ator autorizado em finalidade definida, sem ocultar incerteza. & Registro de uso, decisão, responsável, artefatos considerados, restrições e resultado esperado. \\
Avaliação reflexiva & Execução comparada às referências aplicáveis, sem confundir ausência de evidência com conformidade. & Referências congeladas, evidências consideradas, avaliador/revisão, resultado, divergências, limitações, recomendação e impacto. \\
Retroalimentação & Efeito observado retorna ao ciclo e pode confirmar, revisar ou invalidar conhecimento, política ou ação. & Nova observação relacionada à decisão, comparação e evento de revisão/reprocessamento. \\
\bottomrule
\end{longtable}

\section{Arquitetura funcional e fronteiras de autoridade}

\subsection{Participantes}
\begin{center}
\begin{tikzpicture}[
  node distance=10mm and 14mm,
  comp/.style={draw=sisterblue, rounded corners, fill=sisterlightblue, minimum width=3.7cm, minimum height=12mm, align=center, font=\small\bfseries},
  ext/.style={draw=sisterdarkgray, rounded corners, fill=sistergray, minimum width=3.7cm, minimum height=12mm, align=center, font=\small\bfseries},
  arr/.style={-{Latex[length=2.2mm]}, thick, draw=sisterdarkgray}
]
\node[ext] (human) {Pessoa / automação / CI-CD};
\node[comp, below=of human] (ctl) {\texttt{sisterctl}\\interface de controle};
\node[comp, below=of ctl] (d) {\texttt{sisterd}\\plano de controle persistente};
\node[ext, below left=of d] (sub1) {Subsistema A\\autoridade do domínio};
\node[ext, below right=of d] (sub2) {Subsistema B\\autoridade do domínio};
\node[comp, right=of d] (sge) {SGE-SisTer\\avaliação e evidência};
\node[comp, left=of d] (store) {Persistência\\registros e proveniência};
\draw[arr] (human)--(ctl); \draw[arr] (ctl)--(d);
\draw[arr] (d)--(sub1); \draw[arr] (d)--(sub2);
\draw[arr] (d)--(sge); \draw[arr] (d)--(store);
\draw[arr, dashed] (sub1) to[bend right=15] node[below, font=\scriptsize]{relação contratada} (sub2);
\end{tikzpicture}
\end{center}

\subsection{Matriz de autoridade}
\begin{longtable}{P{3.3cm}P{5.7cm}P{5.8cm}}
\toprule
\textbf{Âmbito} & \textbf{Detém autoridade sobre} & \textbf{Não deve absorver} \\
\midrule
\endhead
Subsistema & Dados nativos, regras especializadas, semântica interna, operação e interface do domínio. & Governança transversal, contratos globais ou dados de autoridade de outros domínios. \\
Relacionamento bilateral & Mapeamento semântico, tradução de identificadores, restrições, correlação e capacidades negociadas. & Regras internas completas de qualquer participante ou decisões globais de maturidade. \\
\texttt{sisterd} & Registro federativo, contratos ativos, autorização transversal, execução, proveniência, estado, conhecimento transversal, avaliações reflexivas, ações corretivas autorizadas e auditoria. & Regras especializadas, bancos de autoridade de todos os domínios, interpretação científica sem método e responsável ou alteração autônoma de referências estruturais. \\
\texttt{sisterctl} & Comando, inspeção, diagnóstico, explicação de divergências, aprovação ou aplicação autorizada de correções e exportação. & Estado próprio divergente, acesso direto à persistência no modo operacional, avaliação oculta no cliente ou regras que pertencem ao serviço. \\
SGE-SisTer & Diagnóstico, atestação, comparação entre evidências e referências, maturidade, promoção, regressão, divergência e publicação. & Execução de regras de domínio, definição unilateral do propósito ou promoção baseada apenas em observação humana. \\
Coordenação científica & Finalidade, método, critérios de avaliação, validade e limites de artefatos científicos. & Implementação técnica, operação cotidiana ou alteração direta de contratos. \\
\bottomrule
\end{longtable}

\section{Catálogo de objetivos sistêmicos}

Os objetivos abaixo vinculam a finalidade científica às funções, processos e capacidades.

\begin{longtable}{P{1.5cm}P{5.1cm}P{8.2cm}}
\toprule
\textbf{ID} & \textbf{Objetivo} & \textbf{Resultado verificável} \\
\midrule
\endhead
O-01 & Governar o ecossistema federativo & Participantes e relacionamentos são reconhecidos, autorizados, versionados, auditáveis e reversíveis. \\
O-02 & Produzir informação integrada rastreável & Dados contratados são validados, contextualizados e relacionados sem perda de origem, autoridade ou incerteza. \\
O-03 & Registrar conhecimento contextualizado e seu uso & Artefatos tipificados são sustentados por informação, método, responsabilidade, limitações, revisão e finalidade. \\
O-04 & Sustentar operação e evolução resilientes & Capacidades são diagnosticáveis, observáveis, recuperáveis e avaliadas por evidências e maturidade reversível. \\
O-05 & Sustentar reflexividade operacional & Execuções relevantes são comparadas às referências vigentes, divergências são explicáveis e respostas corretivas são autorizadas, auditáveis e reavaliadas. \\
\bottomrule
\end{longtable}

\section{Catálogo de funções sistêmicas}

Os identificadores \texttt{F-01} a \texttt{F-11} são estáveis. Uma função pode ser realizada por vários componentes, mas deve possuir um responsável final pela sua definição e critérios de aceitação.

\begin{landscape}
\small
\begin{longtable}{P{1.35cm}P{3.4cm}P{6.2cm}P{4.8cm}P{3.4cm}P{2.7cm}}
\toprule
\textbf{ID} & \textbf{Função} & \textbf{Responsabilidade sistêmica} & \textbf{Saída/evidência mínima} & \textbf{Executor principal} & \textbf{Baseline apps} \\
\midrule
\endhead
F-01 & Reconhecimento federativo & Reconhecer identidade, versão, domínio, responsável, endpoints, capacidades, dependências e perfil de cada participante. & \texttt{SystemRegistration}, manifesto validado, catálogo consultável e histórico. & \texttt{sisterd}; declaração pelo subsistema; inspeção pelo \texttt{sisterctl}. & \statusparcial \\
F-02 & Governança dos relacionamentos & Propor, validar, aceitar, ativar, suspender, revogar e versionar acordos e capacidades autorizadas. & \texttt{IntegrationAgreement}, recibos, digest, revisão, estado e histórico. & \texttt{sisterd} e participantes; comando pelo \texttt{sisterctl}. & \statusparcial \\
F-03 & Identidade, autorização e confiança & Autenticar pessoas e sistemas e autorizar por papel, capacidade, finalidade, escopo e risco. & Sessão ou credencial, decisão de autorização, ator e trilha de operação. & \texttt{sisterd}; identidade de domínio nos subsistemas. & \statusparcial \\
F-04 & Intercâmbio e ingestão & Receber, localizar ou encaminhar dados por API, eventos, pacotes, sincronização e referências. & Entrada identificada, recibo, estado da transferência, limites e erros. & Adaptadores e subsistemas coordenados pelo \texttt{sisterd}. & \statusparcial \\
F-05 & Validação e conformidade & Validar schema, versão, integridade, contrato, origem, finalidade e compatibilidade. & Resultado por verificação, bloqueios, advertências e artefatos validados. & \texttt{sisterd}, contratos e SGE. & \statusparcial \\
F-06 & Evidência e proveniência & Preservar evidências e relacionar origem, execução, transformação, contrato, revisão e derivado. & Evidência com integridade verificável e registro causal consultável. & \texttt{sisterd} e repositórios especializados. & \statusparcial \\
F-07 & Contextualização e correlação & Relacionar dados a entidades, projetos, atividades, tempo, espaço e outras fontes, explicitando pendências. & \texttt{IntegratedInformation}, correlações, qualidade, incertezas e método. & Integrações e motores especializados; governança pelo \texttt{sisterd}. & \statusausente \\
F-08 & Síntese e produção de conhecimento & Registrar interpretações sustentadas por informações, método, confiança, limitações e autoria. & \texttt{KnowledgeArtifact} revisável, publicável e invalidável. & Coordenação científica, serviços de domínio e \texttt{sisterd}. & \statusausente \\
F-09 & Convergência e exposição & Apresentar visão transversal e disponibilizar resultados segundo papel, finalidade e escopo. & API/CLI/interface com estado, proveniência, restrições e versão. & \texttt{sisterd}, \texttt{sisterctl} e interfaces. & \statusparcial \\
F-10 & Diagnóstico e governança operacional & Expor saúde, prontidão, dependências, maturidade, degradações e capacidade de recuperação. & Diagnóstico estruturado, atestação, logs seguros, promoção/regressão e indicadores. & Subsistemas, \texttt{sisterd}, \texttt{sisterctl} e SGE. & \statusparcial \\
F-11 & Reflexividade operacional & Aplicar \texttt{ReflexivityProfile}s, comparar execuções relevantes às referências, separar resultado, gate e ação, explicar divergências e governar respostas dentro da autoridade concedida. & \texttt{ReflexivityProfile}, \texttt{ReferenceSnapshot}, \texttt{OperationalAssessment}, efeito de gate, recomendação, decisão e verificação da correção. & \texttt{sisterd} persiste e coordena; SGE avalia; \texttt{sisterctl} explica e opera; owners definem referências e autorizam. & \statusausente \\
\bottomrule
\end{longtable}
\end{landscape}

Na linha de base, F-03 cobre autenticação humana e autorização por papéis; identidade de sistemas, finalidade e escopo permanecem ausentes. F-06 cobre custódia e consulta de evidências; proveniência causal não foi observada. F-11 não está materializada: diagnósticos e artefatos de maturidade existem como insumos, mas não há referência congelada, comparação reproduzível, resultado reflexivo ou ciclo governado de correção.

\section{Processos finalísticos}

Cada processo deve possuir estado persistente ou evidência suficiente para reconstruir sua execução. Os fluxos abaixo são normativos para o modelo-alvo.

\subsection{P-01 -- Entrada de um subsistema}
\textbf{Objetivo:} reconhecer um participante sem transferir ao núcleo sua autoridade de domínio.

\textbf{Fluxo mínimo:}
\begin{lstlisting}
manifesto apresentado
-> identidade e responsavel validados
-> capacidades e endpoints verificados
-> diagnostico executado
-> registro operacional criado
-> perfil de maturidade associado por evidencia
\end{lstlisting}

\textbf{Saídas:} \texttt{SystemRegistration}, capacidades, responsável, perfil e atestação.  
\textbf{Falhas:} identidade duplicada, schema inválido, raiz não resolvível, owner ausente, diagnóstico inexistente.  
\textbf{Pronto quando:} o sistema pode ser descoberto, inspecionado, diagnosticado e retirado sem alterar código do núcleo.

\subsection{P-02 -- Estabelecimento de um relacionamento}
\textbf{Objetivo:} autorizar uma relação por finalidade e capacidade, não apenas por compatibilidade técnica.

\begin{lstlisting}
proposta
-> validacao dos participantes e capacidades
-> negociacao de finalidade, escopo e restricoes
-> aceite dos owners
-> digest e revisao registrados
-> ativacao
-> observacao, suspensao, revogacao ou substituicao
\end{lstlisting}

\textbf{Saídas:} \texttt{IntegrationAgreement} ativo e recibos dos participantes.  
\textbf{Falhas:} capacidade inexistente, semântica incompatível, finalidade não autorizada, aceite incompleto.  
\textbf{Pronto quando:} nenhuma execução pode ocorrer sem acordo ativo e reproduzível.

\subsection{P-03 -- Recepção de um dado}
\textbf{Objetivo:} registrar uma entrada contratada preservando original, identidade e escopo.

\begin{lstlisting}
solicitacao ou evento
-> autenticacao e autorizacao
-> criacao de IntegrationRun
-> validacao de contrato/schema/integridade
-> registro do original ou referencia imutavel
-> recibo de ingestao
\end{lstlisting}

\textbf{Saídas:} entrada registrada, digest, \texttt{IntegrationRun} e evidência de validação.  
\textbf{Falhas:} origem desconhecida, contrato inativo, schema inválido, excedente de tamanho, indisponibilidade.  
\textbf{Pronto quando:} a entrada pode ser localizada, reproduzida e recusada sem ambiguidade.

\subsection{P-04 -- Produção de informação integrada}
\textbf{Objetivo:} contextualizar e correlacionar dados sem apagar discrepâncias ou incertezas.

\begin{lstlisting}
entrada validada
-> resolucao de entidades e identificadores
-> associacao temporal e espacial
-> correlacao com outras ofertas autorizadas
-> registro de pendencias, qualidade e inconsistencias
-> producao de IntegratedInformation
\end{lstlisting}

\textbf{Saídas:} \texttt{IntegratedInformation}, relações, método, pendências e proveniência.  
\textbf{Falhas:} entidade não resolvida, conflito semântico, fonte indisponível, transformação não reproduzível.  
\textbf{Pronto quando:} toda relação pode ser confirmada, contestada e reprocessada.

\subsection{P-05 -- Produção de conhecimento}
\textbf{Objetivo:} registrar afirmação ou interpretação tipificada, sustentada por informação, método e responsabilidade.

\begin{lstlisting}
informacoes integradas elegiveis
-> selecao do tipo de artefato
-> aplicacao de metodo ou regra declarada
-> formulacao da afirmacao
-> associacao de evidencias e limitacoes
-> atribuicao de autoria, confianca e finalidade
-> revisao pela autoridade aplicavel ao tipo
-> publicacao, invalidacao ou substituicao
\end{lstlisting}

\textbf{Tipos mínimos:} \texttt{scientific}, \texttt{operational}, \texttt{administrative}, \texttt{methodological} e \texttt{governance}.  
\textbf{Saídas:} \texttt{KnowledgeArtifact} tipificado e versionado.  
\textbf{Falhas:} tipo indefinido, método ausente, evidência insuficiente, conflito não tratado ou autoridade revisora inexistente.  
\textbf{Pronto quando:} o artefato pode ser revisado pela autoridade adequada, invalidado e remontado às informações que o sustentam.

\subsection{P-06 -- Uso e retroalimentação}
\textbf{Objetivo:} relacionar conhecimento a decisões e retornar seus efeitos ao ciclo de observação.

\textbf{Saídas:} registro de uso, ator, finalidade, decisão, resultado esperado e nova observação relacionada.  
\textbf{Falhas:} uso fora de escopo, artefato inválido, incerteza omitida ou decisão sem responsável.  
\textbf{Pronto quando:} é possível saber que conhecimento apoiou qual decisão e qual efeito foi observado.

\subsection{P-07 -- Suspensão, invalidação e reconfiguração}
\textbf{Objetivo:} preservar coerência quando contratos, dados, métodos, dependências ou evidências mudam.

\begin{lstlisting}
evento de mudanca ou falha
-> identificacao de execucoes e derivados afetados
-> suspensao ou marcacao de invalidade
-> notificacao dos responsaveis
-> reprocessamento ou substituicao
-> preservacao do historico
\end{lstlisting}

\textbf{Saídas:} evento de impacto, lista de derivados, decisão, reprocessamento e histórico.  
\textbf{Pronto quando:} mudanças relevantes produzem impacto, suspensão, invalidação ou reprocessamento explícitos.

\subsection{P-08 -- Avaliação reflexiva e ação corretiva}
\textbf{Objetivo:} verificar se uma execução relevante permaneceu coerente com as referências que a autorizaram e governar a resposta sem ultrapassar o perfil de reflexividade aplicável.

\begin{lstlisting}
execucao relevante concluida ou interrompida
-> selecao e validacao do ReflexivityProfile
-> congelamento das referencias aplicaveis
-> coleta e validacao das evidencias
-> comparacao por avaliador versionado
-> resultado: confirmed | divergent | inconclusive | not_applicable
-> determinacao independente de severidade e efeito de gate
-> resposta: none | collect_evidence | recommend | request_approval
            | cancel | retry | isolate | rollback | suspend | reprocess
-> decisao conforme autoridade A0--A5
-> aplicacao reversivel quando autorizada
-> nova execucao e verificacao de efetividade
\end{lstlisting}

\textbf{Saídas:} \texttt{ReflexivityProfile}, \texttt{ReferenceSnapshot}, \texttt{OperationalAssessment}, divergências tipificadas, efeito de gate, recomendação, decisão, \texttt{CorrectiveAction} e avaliação posterior.  
\textbf{Falhas:} perfil ausente ou incompatível, referência ausente ou mutável, evidência insuficiente, avaliador sem versão, confusão entre resultado/gate/ação, conflito de autoridade, correção não reversível ou ação sem verificação posterior.  
\textbf{Pronto quando:} a conclusão pode ser reproduzida, explicada e contestada; o efeito aplicado corresponde ao perfil; e nenhuma ação além do nível autorizado é executada automaticamente.

P-08 é transversal: pode ser acionado ao final de P-01 a P-07, por evento crítico, mudança de referência ou avaliação periódica. Ele não substitui os processos finalísticos; observa seu desempenho e fecha o ciclo entre evidência, decisão e reconfiguração. O primeiro uso normativo é Nexo--Compras em D2--D3, A1 e \texttt{shadow}.

\section{Especificação dos componentes}

\subsection{\texttt{sisterd}: plano de controle persistente}
\textbf{Por que existe:} manter o estado governado do ecossistema e coordenar funções transversais que não pertencem a um único domínio.

\textbf{Deve:}
\begin{itemize}
  \item manter registros de sistemas, capacidades, acordos, execuções, evidências, proveniência e conhecimento;
  \item autenticar e autorizar operações por capacidade, finalidade e escopo;
  \item coordenar integrações sem incorporar regras internas dos domínios;
  \item expor APIs estáveis, saúde, prontidão, auditoria e evidências consumíveis pelo SGE;
  \item suportar suspensão, invalidação, reprocessamento e impacto sobre derivados;
  \item registrar e validar \texttt{ReflexivityProfile}s, congelar referências aplicáveis, persistir avaliações reflexivas e coordenar ações corretivas autorizadas e sua verificação posterior.
\end{itemize}

\textbf{Não deve:}
\begin{itemize}
  \item ser banco de autoridade de todos os subsistemas;
  \item produzir inferência científica sem método e responsável explícitos;
  \item ocultar indisponibilidade por fallback indistinguível de dado real;
  \item decidir sozinho promoção de maturidade;
  \item concentrar novas responsabilidades no arquivo monolítico existente;
  \item alterar autonomamente propósito, contrato, política, critério científico ou código em resposta a uma divergência.
\end{itemize}

\subsection{\texttt{sisterctl}: interface operacional e administrativa}
\textbf{Por que existe:} permitir operação humana e automação sem acesso direto à persistência ou a interfaces privadas no modo operacional normal. Comandos de \textit{bootstrap}, migração, manutenção ou recuperação podem acessar mecanismos internos quando forem separados, autorizados e auditáveis.

\textbf{Todo comando permanente deve possuir:}
\begin{itemize}
  \item sintaxe documentada e parâmetros explícitos;
  \item saída humana e saída estruturada (\texttt{--json});
  \item códigos de retorno estáveis;
  \item modo não interativo para CI/CD;
  \item confirmação para ação destrutiva;
  \item idempotência quando aplicável;
  \item ator, ambiente, alvo, operação e evidência relacionados.
\end{itemize}

\textbf{Comandos-alvo:}
\begin{lstlisting}
sisterctl system list|show|register|retire
sisterctl capability list|show
sisterctl agreement propose|validate|activate|suspend|revoke
sisterctl integration run|status|result|reprocess
sisterctl knowledge list|show|review|invalidate|publish
sisterctl audit show|export
sisterctl doctor
sisterctl maturity evaluate|publish
sisterctl reflexivity profile list|show|validate
sisterctl assessment list|show|explain
sisterctl correction propose|approve|apply|rollback|verify
\end{lstlisting}

\subsection{Subsistemas}
Cada subsistema declara identidade, capacidades, contratos, diagnóstico, dados de autoridade, política de falha, referências de avaliação aplicáveis e responsável. Ele executa regras especializadas e fornece evidências suficientes para que o SisTer governe e avalie a relação, sem entregar ao núcleo detalhes internos desnecessários.

\subsection{Adaptadores de integração}
Adaptadores traduzem protocolos e contratos técnicos. Eles devem ser pequenos, substituíveis, testáveis e incapazes de redefinir a semântica dos domínios. Toda transformação relevante deve possuir identificador, versão, entrada, saída, erro e evidência.

\subsection{SGE-SisTer}
O SGE avalia componentes, integrações, processos e evidências. Deve operar com motores explicitamente identificados, registrar divergências e permitir estados como \texttt{PASS}, \texttt{FAIL}, \texttt{WARN}, \texttt{SHADOW}, \texttt{NOT\_APPLICABLE}, \texttt{SKIP\_JUSTIFIED}, \texttt{UNRESOLVED}, \texttt{BLOCKED} e \texttt{ERROR}. Na reflexividade operacional, valida o \texttt{ReflexivityProfile}, compara evidências a \texttt{ReferenceSnapshot}s, produz \texttt{OperationalAssessment}s explicáveis e publica separadamente resultado, severidade e efeito de gate. O SGE não executa ação corretiva por inferência nem redefine unilateralmente as referências avaliadas.

\subsection{CI/CD e publicação}
A esteira não cria maturidade; ela executa gates e publica decisões sustentadas por evidências. Nenhuma capacidade pode avançar sem perfil aplicável, evidência estruturada ou justificativa válida para \texttt{SKIP}. Quando P-08 for aplicável, a esteira também deve publicar a avaliação reflexiva e impedir promoção diante de divergência bloqueante ou inconclusão não tratada.

\section{Linha de base verificada do código}

\subsection{Inventário inspecionado}
\begin{longtable}{P{5.0cm}P{1.6cm}P{7.7cm}}
\toprule
\textbf{Arquivo/módulo} & \textbf{Linhas} & \textbf{Responsabilidade observada} \\
\midrule
\endhead
\texttt{apps/sisterd/}\newline\texttt{main.cpp} & 2.267 & Servidor HTTP próprio, configuração, concorrência, parsing, autenticação de rotas, proxy reverso, APIs, arquivos estáticos, logs e despacho principal. \\
\texttt{apps/sisterd/}\newline\texttt{auth.cpp/.hpp} & 604 & Usuários, papéis, senha com salt/hash, sessões persistidas, bootstrap, login, CRUD administrativo e logout. \\
\texttt{apps/sisterd/}\newline\texttt{db.cpp/.hpp} & 193 & Conexão opcional PostgreSQL e consultas de leitura para sistemas, contratos, evidências e diagnósticos; stub quando compilado sem libpq. \\
\texttt{apps/sisterd/}\newline\texttt{studio\_client.cpp/.hpp} & 206 & Cliente TLS para capacidades e saúde do Studio, verificação de CA/host e token por arquivo. \\
\texttt{apps/sisterd/}\newline\texttt{sister\_campo\_client.}\allowbreak\texttt{cpp/.hpp} & 106 & Cliente HTTP local, validação mínima de saúde/capacidade e estado integrado/degradado. \\
\texttt{apps/sisterd/api/}\newline\texttt{maturity\_routes.cpp/.hpp} & 169 & Exposição administrativa de artefatos de maturidade publicados, com limite de tamanho e verificação de identificador de schema. \\
\texttt{apps/sisterctl/}\newline\texttt{main.cpp} & 158 & Validação de manifesto, verificação/migração de banco por scripts e importação interativa de usuário. \\
\bottomrule
\end{longtable}

\subsection{Capacidades efetivamente observadas no \texttt{sisterd}}
\begin{itemize}
  \item configuração por variáveis de ambiente, limites numéricos e defaults diferenciados por ambiente;
  \item servidor HTTP com fila limitada, conjunto de workers, timeouts, encerramento por sinal e identificador de requisição;
  \item autenticação por sessão, cookie, bootstrap, login, logout, rate limiting de falhas e gestão administrativa de usuários;
  \item autorização atualmente centrada em papéis, principalmente \texttt{admin}, com catálogo de capacidades derivado do papel;
  \item proteção de origem para métodos inseguros, cookie seguro configurável, HSTS opcional e controle de arquivos estáticos;
  \item APIs de saúde, sistemas, contratos, evidências, diagnósticos e descritores de integrações;
  \item proxy autenticado para Nexo e Clima; cliente TLS específico para Studio; verificação local específica para SisTer-Campo;
  \item leitura PostgreSQL de registros selecionados e operação degradada/fallback quando o banco não responde ou libpq não está presente;
  \item exposição de artefatos administrativos de maturidade, catálogo, componentes, qualidade, última execução e histórico.
\end{itemize}

\subsection{Endpoints observados}
\begin{landscape}
\small
\begin{longtable}{P{4.5cm}P{3.0cm}P{3.2cm}P{8.0cm}}
\toprule
\textbf{Endpoint} & \textbf{Método} & \textbf{Controle} & \textbf{Comportamento observado} \\
\midrule
\endhead
\texttt{/api/health} & GET/HEAD & público & Sinal de vida e diagnóstico básico; informa versão e conexão com banco, mas não constitui prontidão completa. \\
\texttt{/api/auth/bootstrap} & GET & público controlado & Informa disponibilidade de bootstrap administrativo. \\
\texttt{/api/auth/register} & POST & bootstrap & Registra primeiro administrador e cria sessão. \\
\texttt{/api/auth/login} & POST & público com limite & Autentica e cria cookie/sessão. \\
\texttt{/api/auth/logout} & POST & autenticado & Revoga sessão e limpa cookie. \\
\texttt{/api/me}, \texttt{/api/me/capabilities} & GET & autenticado & Usuário atual e capacidades derivadas de papel. \\
\texttt{/api/admin/users[/id]} & \shortstack[l]{GET/POST/PUT/\\PATCH/DELETE} & administrador & Gestão de contas. \\
\texttt{/api/systems} & GET & autenticado & Consulta PostgreSQL ou catálogo fallback. \\
\texttt{/api/contracts} & GET & administrador & Consulta PostgreSQL ou catálogo fallback. \\
\texttt{/api/evidence} & GET & administrador & Consulta de evidências; retorna lista vazia como fallback. \\
\texttt{/api/diagnostics} & GET & administrador & Consulta PostgreSQL ou diagnóstico fallback. \\
\texttt{/api/integrations/*} & GET & varia por integração & Descritor fixo ou resultado de cliente/adaptador específico. \\
\texttt{/integrations/nexo/*} & vários/proxy & autenticado & Proxy reverso local com identidade encaminhada. \\
\texttt{/integrations/clima/*} & vários/proxy & autenticado & Proxy HTTP/WebSocket local com identidade encaminhada. \\
\texttt{/api/admin/maturity/*} & GET/HEAD & administrador & Lê e valida artefatos JSON publicados em diretório configurado. \\
\bottomrule
\end{longtable}
\end{landscape}

\subsection{Capacidades observadas no \texttt{sisterctl}}
\begin{tabularx}{\textwidth}{P{4.7cm}Y}
\toprule
\textbf{Comando atual} & \textbf{Comportamento} \\
\midrule
\texttt{validate-manifest <arquivo>} & Carrega JSON, usa contrato comum e retorna sucesso/falha de validação. \\
\texttt{db-check} & Delega a \texttt{scripts/dev/db\_check.sh}. \\
\texttt{db-migrate [arquivo]} & Delega a \texttt{scripts/dev/db\_migrate.sh}, com quoting do argumento. \\
\texttt{auth-import-user ...} & Solicita senha sem eco e escreve por meio de \texttt{AuthStore}. \\
\bottomrule
\end{tabularx}

\subsection{O que o snapshot ainda não materializa}
Não foram observados, no código inspecionado:
\begin{itemize}
  \item API de escrita e ciclo de vida para registro federativo e acordos;
  \item modelo persistente e API de \texttt{IntegrationRun};
  \item registro causal de proveniência entre original, transformação e derivado;
  \item operação genérica de validação por schema/contrato nas integrações;
  \item \texttt{IntegratedInformation} como artefato governado;
  \item \texttt{KnowledgeArtifact} e seus estados de revisão, publicação e invalidação;
  \item consulta de impacto e reprocessamento de derivados;
  \item trilha de auditoria de domínio exportável;
  \item cliente \texttt{sisterctl} para as APIs do \texttt{sisterd}, saída \texttt{--json} e modo não interativo abrangente;
  \item autorização baseada diretamente em capacidades, finalidade e escopo; no snapshot, o papel administrativo ainda concentra decisões;
  \item diagnóstico federado acionável por \texttt{sisterctl doctor};
  \item referência congelada, avaliação reflexiva reproduzível, explicação de divergências e ciclo autorizado de ação corretiva.
\end{itemize}

\begin{alerta}[Diagnóstico arquitetural do snapshot]
O código já materializa uma base de operação e segurança considerável, mas o arquivo \texttt{main.cpp} concentra múltiplas camadas e responsabilidades. A próxima evolução deve preservar o comportamento atual, extrair fronteiras mínimas e implementar processos finalísticos por fatias verticais.
\end{alerta}

\section{Matriz atual versus alvo}

Os estados de implementação registram presença observada no snapshot e não equivalem a maturidade. A conformidade compara o comportamento observado com o modelo-alvo.

\begin{landscape}
\scriptsize
\begin{longtable}{P{3.2cm}P{3.0cm}P{2.5cm}P{5.1cm}P{5.6cm}}
\toprule
\textbf{Capacidade} & \textbf{Implementação} & \textbf{Conformidade} & \textbf{Código/evidência atual} & \textbf{Lacuna para o objetivo} \\
\midrule
\endhead
Servidor e configuração & \statusimpl & Parcial & HTTP, workers, fila, timeouts, configuração, logs e arquivos estáticos. & Extrair responsabilidades, caracterizar protocolo, decidir manutenção do servidor HTTP e implementar prontidão estruturada. \\
Identidade de usuários & \statusimpl & Parcial & Bootstrap, login, sessões, CRUD, papéis e logout. & Capacidade, finalidade, escopo, identidade de sistemas e credenciais adequadas à produção. \\
Registro de sistemas & \statusparcial & Parcial & Consulta PostgreSQL, fallback e validação de manifesto via CLI. & Registro, versão, owner, capacidades, estado operacional, perfil de maturidade, retirada e histórico. \\
Contratos & \statusparcial & Parcial & Consulta de contratos e descritores fixos de integrações. & Ciclo completo, aceite bilateral, digest, suspensão, revogação e compatibilidade. \\
Integrações & \statusparcial & Parcial & Proxies e clientes específicos para Nexo, Clima, Studio e Campo. & Adaptadores contratados, execução identificada, recibos, idempotência, timeout e política de retry. \\
Validação & \statusparcial & Parcial & Manifesto, identificadores de schema em maturidade e verificações específicas. & Validação por contrato/schema, erros estruturados, versão e evidência por check. \\
Evidências & \statusparcial & Parcial & Consulta de evidências e artefatos de maturidade publicados. & Escrita, integridade, retenção, classificação e vínculo causal. \\
Proveniência causal & \statusausente & Não conforme & Não observada no snapshot. & Relação origem--execução--transformação--derivado, consulta de impacto e reprocessamento. \\
Diagnóstico & \statusparcial & Parcial & \texttt{/api/health}, consultas/fallback e maturidade exposta. & Separar \textit{liveness}, \textit{readiness} e degradação; implementar \texttt{doctor} federado. \\
Reflexividade operacional & \statusespec & Não conforme & Logs, diagnósticos e artefatos de maturidade fornecem insumos dispersos, mas não constituem avaliação reflexiva. & Congelar referências; comparar evidências; produzir \texttt{OperationalAssessment}; explicar divergências; autorizar, aplicar e verificar correções. \\
\texttt{IntegrationRun} & \statusespec & Não conforme & Estruturas citadas nos documentos, não observadas no snapshot. & Entidade, estado de execução, validade, relações entre execuções, persistência, API e evidência. \\
Informação integrada & \statusespec & Não conforme & Processo Nexo--Compras descrito fora do snapshot. & Artefato, critérios, correlação, qualidade, pendências e proveniência. \\
\texttt{KnowledgeArtifact} & \statusespec & Não conforme & Modelo conceitual nos documentos. & Tipos, estados independentes de revisão/publicação/validade, autoridade revisora e impacto. \\
CLI governada & \statusparcial & Parcial & Quatro operações locais, incluindo migração e importação. & Cliente HTTP, comandos funcionais, JSON, não interativo, códigos estáveis e separação de comandos de manutenção. \\
\bottomrule
\end{longtable}
\end{landscape}

\section{Arquitetura de implementação proposta}

\subsection{Dependências entre camadas}
A estrutura abaixo define a direção de dependência do modelo-alvo.

\begin{center}
\begin{tikzpicture}[
  node distance=10mm and 8mm,
  layer/.style={draw=sisterblue, rounded corners, text width=7.8cm, minimum height=11mm, align=center, font=\small\bfseries},
  infra/.style={draw=sisterdarkgray, rounded corners, fill=sistergray, text width=4.6cm, minimum height=18mm, align=center, font=\small\bfseries},
  arr/.style={-{Latex[length=2.2mm]}, thick, draw=sisterdarkgray},
  impl/.style={-{Latex[length=2.2mm]}, dashed, thick, draw=sisterdarkgray}
]
\node[layer, fill=sisterlightblue] (api) {Interfaces: HTTP, CLI, eventos e pacotes};
\node[layer, fill=sistergray, below=of api] (app) {Aplicação\\casos de uso, comandos, consultas, transações e portas};
\node[layer, fill=sisterlightblue, below=of app] (dom) {Domínio\\entidades, estados, invariantes, políticas e eventos};
\node[infra, right=of app] (infra) {Infraestrutura\\PostgreSQL, arquivos, rede, TLS, filas e adaptadores};
\draw[arr] (api)--(app);
\draw[arr] (app)--(dom);
\draw[impl] (infra)--node[above, font=\scriptsize]{implementa portas} (app);
\draw[impl] (infra.south west)--(dom.east);
\end{tikzpicture}
\end{center}

\begin{decisao}[Regra de dependência]
O domínio não depende de HTTP, PostgreSQL, arquivos, terminal ou build. A aplicação depende do domínio e define portas. Interfaces acionam casos de uso. A infraestrutura implementa portas e depende das abstrações internas.
\end{decisao}

\subsection{Estratégia de extração incremental}
A decomposição deve preservar o comportamento atual e avançar por fatias verticais:
\begin{enumerate}
  \item criar testes de caracterização das rotas, autenticação, códigos e comandos existentes;
  \item registrar ADR sobre manutenção ou substituição do servidor HTTP próprio;
  \item extrair transporte, roteamento e composição sem alterar contratos externos;
  \item criar portas mínimas de persistência, relógio, identidade e integração;
  \item implementar uma fatia Nexo--Compras com \texttt{IntegrationRun}, proveniência e informação integrada;
  \item generalizar módulos somente após comprovação por teste e evidência.
\end{enumerate}

\subsection{Decisão sobre o servidor HTTP}
O servidor HTTP próprio concentra parsing, sockets, concorrência, proxy, WebSocket, cookies, origem, limites e arquivos estáticos. Sua continuidade exige ADR com alternativa considerada, responsabilidade, perfil de exposição, testes de protocolo, entradas malformadas, contenção de exceções, limites de recursos, atualização e segurança.

\begin{principio}[Princípio de transporte restrito]
\textbf{O SisTer não pretende reproduzir um servidor HTTP generalista. Quando implementar ou utilizar capacidades de transporte, deverá herdar o conhecimento acumulado sobre suas ameaças, limitar sua superfície e demonstrar por evidências que os riscos aplicáveis estão controlados.}
\end{principio}

Por padrão, o \texttt{sisterd} é um plano de controle interno. TLS público, normalização de borda, conexões persistentes, WebSocket, arquivos estáticos, quotas externas e funções de proxy generalista devem ser delegados a componentes especializados. Exceção exige ADR, modelo de ameaças, testes negativos e risco residual explícito.

\subsection{Módulos de domínio propostos}
\begin{tabularx}{\textwidth}{P{4.3cm}Y}
\toprule
\textbf{Módulo} & \textbf{Responsabilidade} \\
\midrule
\texttt{federation} & Sistemas, versões, owners, capacidades, endpoints, estado operacional e perfil de maturidade. \\
\texttt{agreements} & Proposta, aceite, ativação, suspensão, revogação, digest e compatibilidade. \\
\texttt{execution} & \texttt{IntegrationRun}, estado de execução, idempotência, bloqueios, tentativas e relações entre execuções. \\
\texttt{provenance} & Origem, transformação, evidência, derivação, impacto e reprocessamento. \\
\texttt{information} & Contextualização, correlação, qualidade, pendências e \texttt{IntegratedInformation}. \\
\texttt{knowledge} & Tipo, afirmação, método, confiança, revisão, publicação, validade e uso. \\
\texttt{identity} & Pessoas, sistemas, credenciais, papéis, capacidades, finalidade e escopo. \\
\texttt{operations} & \textit{Liveness}, \textit{readiness}, degradação, diagnóstico, auditoria, métricas e SGE. \\
\texttt{reflexivity} & Referências de avaliação, comparação, divergências, recomendações, decisões corretivas, reversibilidade e verificação de efetividade. \\
\bottomrule
\end{tabularx}


\subsection{Diretrizes de materialização da reflexividade em C++}
\label{sec:cpp-reflexividade}

A implementação deve distinguir recursos estáveis da linha de base e experimentos de linguagem. O modelo funcional não pode depender de recurso ainda indisponível na matriz oficial de compiladores do projeto.

\begin{tabularx}{\textwidth}{P{3.0cm}P{5.2cm}Y}
\toprule
\textbf{Faixa} & \textbf{Mecanismos} & \textbf{Uso recomendado} \\
\midrule
C++20 & tipos fortes, \texttt{concepts}, \texttt{constexpr}/\texttt{consteval}, \texttt{std::source\_location}, \texttt{std::jthread}, \texttt{std::stop\_token}, \texttt{std::variant} & Invariantes estruturais, contratos de tipo, máquinas de estado, origem técnica e cancelamento cooperativo. \\
C++23 & \texttt{std::expected}, \texttt{std::stacktrace} & Resultados e falhas tipificados, reconstrução técnica de falhas e integração direta com evidências. \\
Laboratório C++26 & contratos da linguagem, reflexão estática e \texttt{std::execution} & Prototipar geração de metadados/schemas, precondições/pós-condições e composição assíncrona, sem tornar o piloto dependente desses recursos. \\
\bottomrule
\end{tabularx}

\subsubsection{Elementos arquiteturais mínimos}
A implementação inicial deve separar cinco responsabilidades:
\begin{lstlisting}[language=C++]
struct ReflexiveContext;
class EvidenceCollector;
class AssessmentEngine;
struct ReflexivityPolicy;
class CorrectiveActionExecutor;
\end{lstlisting}

\begin{itemize}
  \item \texttt{ReflexiveContext}: identifica execução, referência, política, correlação, ator, modo e origem técnica;
  \item \texttt{EvidenceCollector}: captura evidências estruturadas durante todo o escopo da operação;
  \item \texttt{AssessmentEngine}: compara referência e evidência sem efeitos colaterais, com versão e explicação;
  \item \texttt{ReflexivityPolicy}: materializa profundidade, autoridade, tolerâncias, gate e ações permitidas do perfil;
  \item \texttt{CorrectiveActionExecutor}: executa somente ação previamente autorizada e devolve recibo ou falha tipificada.
\end{itemize}

\subsubsection{Regras de projeto C++}
\begin{enumerate}
  \item identificadores de domínios distintos devem ser tipos distintos, evitando conversão implícita entre \texttt{RunId}, \texttt{AgreementId}, \texttt{SystemId} e equivalentes;
  \item estados centrais devem ser modelados por enumerações ou variantes tipificadas, sem strings livres como fonte de verdade;
  \item precondições e invariantes estáveis devem ser verificadas em compilação quando possível e novamente nas fronteiras não confiáveis;
  \item operações com falha esperada devem preferir resultado tipificado, incluindo código, explicação, evidência, possibilidade de retry e indicação de efeitos já confirmados;
  \item coleta de evidência deve usar RAII para registrar conclusão, falha, cancelamento ou abandono, sem depender de todos os caminhos de retorno lembrarem de emitir logs;
  \item cancelamento de tarefas deve ser cooperativo e observável, evitando término abrupto de threads e perda de estado;
  \item o avaliador deve ser determinístico, reproduzível e separado do executor de ações corretivas;
  \item contratos, reflexão e \texttt{std::execution} do C++26 permanecem em laboratório até haver suporte documentado, testes de portabilidade e ADR de adoção.
\end{enumerate}

\begin{lstlisting}[language=C++]
OperationalAssessment assess(
    const ReflexivityProfile& profile,
    const ReferenceSnapshot& reference,
    const EvidenceBundle& evidence);

std::expected<ActionReceipt, ActionFailure>
execute_corrective_action(
    const CorrectiveAction& action,
    const Authorization& authorization);
\end{lstlisting}

\begin{decisao}[Fronteira de autoridade no código]
O componente que detecta ou classifica a divergência não recebe, por esse fato, autoridade para corrigi-la. Avaliação e execução corretiva devem possuir interfaces, políticas, testes e registros separados.
\end{decisao}

\subsubsection{Estratégia para recursos C++26}
Enquanto contratos nativos não integrarem a linha de base, o projeto pode usar uma fachada própria, por exemplo \texttt{SISTER\_PRECONDITION}, \texttt{SISTER\_POSTCONDITION} e \texttt{SISTER\_INVARIANT}, desde que cada violação produza evidência estruturada. A reflexão estática deve ser experimentada para inventário de campos, serialização, geração de schemas e consistência entre modelo, API e CLI. Nenhum experimento pode redefinir contratos externos ou gerar código de produção sem revisão e testes equivalentes ao caminho manual.

\subsection{Estrutura de diretórios sugerida}
\begin{lstlisting}
apps/sisterd/
  main.cpp                  # composicao e inicializacao
  api/                      # controllers, DTOs, serializacao
  application/              # casos de uso e portas
  domain/
    federation/
    agreements/
    execution/
    provenance/
    information/
    knowledge/
    identity/
    operations/
    reflexivity/
  infrastructure/
    postgres/
    filesystem/
    http/
    integrations/

apps/sisterctl/
  main.cpp                  # parser e composicao
  commands/
    operational/
    maintenance/
  client/
  output/
  config/
\end{lstlisting}


\section{Engenharia de segurança orientada por ameaças}
\label{sec:engenharia-seguranca}

\subsection{Finalidade}
A segurança do SisTer deve ser conduzida como processo permanente de engenharia. O objetivo não é prever toda vulnerabilidade, mas manter um sistema capaz de conhecer ameaças já documentadas, restringir superfície, conter falhas, incorporar descobertas e impedir promoção sem evidência.

A EFE define as obrigações normativas. O registro operacional de ameaças, controles, risco residual e revisão deve ser mantido em documento complementar versionado, denominado \textbf{MAES-SisTer -- Modelo de Ameaças e Estratégia de Segurança}. Mudança de ameaça não altera necessariamente a EFE; mudança de função, fronteira, protocolo ou garantia estrutural altera.

\subsection{Ativos e fronteiras de confiança}
Ativos mínimos protegidos:
\begin{itemize}
  \item identidades humanas e de sistemas, credenciais, sessões e chaves;
  \item contratos, acordos, capacidades e decisões de autorização;
  \item dados de autoridade, execuções, proveniência, informação integrada e conhecimento;
  \item evidências, decisões de maturidade e integridade da cadeia de publicação;
  \item disponibilidade, recuperabilidade e configuração do plano de controle.
\end{itemize}

\begin{center}
\begin{tikzpicture}[
  node distance=7mm,
  b/.style={draw=sisterblue, rounded corners, fill=sisterlightblue, minimum width=8.8cm, minimum height=9mm, align=center, font=\small\bfseries},
  a/.style={-{Latex[length=2.2mm]}, thick, draw=sisterdarkgray}
]
\node[b] (net) {Rede externa ou institucional};
\node[b, below=of net] (gw) {Gateway especializado};
\node[b, below=of gw] (sd) {\texttt{sisterd}: plano de controle interno};
\node[b, below=of sd] (clients) {Clientes internos, persistência, \texttt{sisterctl} e SGE};
\node[b, below=of clients] (subs) {Subsistemas autônomos};
\draw[a] (net)--(gw); \draw[a] (gw)--(sd); \draw[a] (sd)--(clients); \draw[a] (clients)--(subs);
\end{tikzpicture}
\end{center}

Cada travessia deve declarar identidade aceita, protocolo, capacidades, dados permitidos, validação, limites, auditoria, comportamento em falha e owner.

\subsection{Ficha obrigatória de ameaça}
Cada ameaça aplicável deve possuir, no mínimo:
\begin{itemize}
  \item identificador, título, ativo e superfície;
  \item origem, precondições, cenário, impacto e exposição;
  \item controles preventivos, detectivos, de contenção e recuperação;
  \item testes, evidências, risco residual, owner, estado e data de revisão;
  \item fontes normativas, históricas ou internas que sustentam sua inclusão.
\end{itemize}

\subsection{Classes iniciais de ameaças}
\begin{landscape}
\begingroup
\tiny
\renewcommand{\arraystretch}{1.02}
\begin{longtable}{P{2.2cm}P{3.2cm}P{6.5cm}P{7.7cm}}
\toprule
\textbf{ID} & \textbf{Classe} & \textbf{Cenário representativo} & \textbf{Estratégia de engenharia} \\
\midrule
\endhead
TH-HTTP-01 & Parsing e enquadramento & \texttt{Content-Length} inválido, duplicado, ambíguo, overflow ou combinação incompatível de headers. & Parsing estrito; erro controlado; conexão encerrada; corpus negativo; conformidade com RFC; teste de sobrevivência do processo. \\
TH-HTTP-02 & Dessincronização HTTP & Gateway e backend interpretam a mesma mensagem de forma diferente, permitindo smuggling ou splitting. & Normalização no gateway; sintaxe restrita; testes diferenciais; rejeição de framing ambíguo; ausência de tolerância incompatível. \\
TH-HTTP-03 & Exaustão de recursos & Requisição lenta, corpo grande, filas, threads, memória, descritores ou logs esgotados. & Limites absolutos e por origem; deadlines; fila limitada; quotas; backpressure; testes de carga e exaustão. \\
TH-HTTP-04 & Injeção de cabeçalhos & CRLF, nomes inválidos, cabeçalhos internos forjados ou resposta maliciosa do upstream. & Allowlist; sanitização; sobrescrita de headers internos; validação de respostas; logs seguros. \\
TH-PROXY-01 & SSRF e proxy indevido & Entrada controla destino, esquema, porta ou caminho e transforma o sistema em proxy aberto. & Destinos fixos e contratados; egress restrito; resolução validada; nenhuma URL arbitrária; testes de negação. \\
TH-PROXY-02 & Upstream comprometido & Subsistema retorna headers, tamanhos, estados ou fluxos maliciosos. & Todo upstream é não confiável; limites de resposta; parser independente; timeout; circuito de contenção; auditoria. \\
TH-AUTH-01 & Autenticação e sessão & Força bruta, enumeração, sessão roubada, bootstrap indevido ou credencial persistente. & Rate limiting multidimensional; resposta uniforme; cookies seguros; bootstrap offline; expiração e revogação. \\
TH-AUTHZ-01 & Autorização ausente ou excessiva & Rota sem política, capacidade ampla ou decisão reinterpretada pelo handler. & Negação por padrão; política declarada na entrada; menor privilégio; teste \texttt{401}/\texttt{403}; auditoria correlacionada. \\
TH-IDENT-01 & Identidade entre sistemas & Cookie humano, header forjado, token sem audiência ou credencial estática atravessa a fronteira. & Asserção interna assinada, curta, com audiência, capacidade e \texttt{request\_id}; rotação; rejeição de replay quando aplicável. \\
TH-WS-01 & Conexões persistentes & WebSocket ocupa worker, não expira, ignora origem ou permite retenção de recursos. & WebSocket no gateway; quotas; origem permitida; timeout de inatividade e vida máxima; proxy legado proibido em produção. \\
TH-PATH-01 & Caminhos e arquivos & Traversal, symlink, encoding ambíguo, leitura indevida ou exposição de fonte. & Raiz fixa; canonicalização; privilégios mínimos; arquivos estáticos fora do \texttt{sisterd}; testes de caminhos hostis. \\
TH-CXX-01 & Memória e inteiros & Overflow, leitura/escrita fora de limites, uso após liberação, ponteiro nulo ou corrida. & Tipos limitados; validação; RAII; sanitizers; análise estática; fuzzing; revisão especializada. \\
TH-CXX-02 & Exceções e encerramento & Entrada externa ou falha de handler propaga exceção e encerra worker ou processo. & Erros esperados tratados localmente; barreira final por job; isolamento; teste hostil seguido de health. \\
TH-CONF-01 & Configuração insegura & Bind público, ambiente desconhecido, proxy legado ou bootstrap habilitado em produção. & Schema e validação no startup; defaults seguros; falha fechada; unit file explícito; teste da configuração efetiva. \\
TH-SUP-01 & Cadeia de suprimento & Dependência vulnerável, build não reproduzível, artefato alterado ou origem não verificada. & Inventário/SBOM; pinning; verificação de origem e integridade; advisories; atualização e rollback. \\
TH-AUD-01 & Logs e evidências & Segredo em log, injeção de controle, ausência de correlação ou evidência alterável. & Sanitização; schema; retenção; integridade; classificação; separação entre log operacional e auditoria governada. \\
\bottomrule
\end{longtable}
\endgroup
\end{landscape}

\subsection{Invariantes para ameaças conhecidas e novas}
\begin{enumerate}
  \item Toda entrada é não confiável, inclusive a produzida por subsistema autenticado.
  \item Falhas são contidas na menor unidade possível e não atravessam fronteiras sem tradução controlada.
  \item Memória, tempo, conexões, filas, threads, corpos, descritores e logs possuem limites explícitos.
  \item Toda fronteira nega por padrão e toda permissão possui capacidade, escopo e owner.
  \item Credenciais não atravessam fronteiras sem necessidade demonstrada.
  \item Configuração desconhecida ou incompleta falha fechada.
  \item O upstream, o gateway e o cliente podem estar defeituosos ou comprometidos.
  \item Segurança depende de camadas independentes, não de um único controle.
  \item Vulnerabilidade corrigida gera teste de regressão e atualização do modelo de ameaças.
  \item Ameaça nova é associada a classe existente ou cria nova classe versionada.
  \item Perda de controle ou evidência pode regredir maturidade, suspender capacidade ou revogar integração.
  \item Capacidade não necessária deve ser removida ou mantida fora da superfície operacional.
\end{enumerate}

\subsection{Ciclo de incorporação de ameaças}
\begin{lstlisting}
descoberta, incidente, advisory ou nova exposicao
        -> avaliacao de aplicabilidade
        -> ameaca existente ou novo TH-ID
        -> risco e fronteira afetada
        -> controle e owner
        -> teste reproduzivel
        -> implementacao e evidencia
        -> risco residual
        -> avaliacao SGE
        -> promocao, manutencao, regressao ou bloqueio
\end{lstlisting}

O acompanhamento deve considerar RFCs e erratas dos protocolos utilizados, advisories dos componentes e dependências, CVE/CWE aplicáveis, sistema operacional, resultados de fuzzing e sanitizers, incidentes internos e vulnerabilidades históricas de servidores e proxies maduros.

\subsection{Verificação proporcional à exposição}
Conforme risco e superfície, o conjunto de evidências deve combinar:
\begin{itemize}
  \item testes unitários, negativos, de contrato, integração e ponta a ponta;
  \item corpus de entradas hostis e testes diferenciais de protocolo;
  \item fuzzing de parsers e fronteiras de serialização;
  \item AddressSanitizer, UndefinedBehaviorSanitizer e ThreadSanitizer quando aplicável;
  \item análise estática, warnings máximos e revisão de dependências;
  \item testes de carga, requisição lenta, exaustão, upstream hostil e recuperação;
  \item verificação de configuração produtiva, sandbox, rollback, backup e restauração.
\end{itemize}

\subsection{Rastreabilidade de segurança}
\begin{principio}[Cadeia de segurança]
\centering
\textbf{Ameaça $\rightarrow$ requisito $\rightarrow$ controle $\rightarrow$ implementação $\rightarrow$ teste $\rightarrow$ evidência $\rightarrow$ risco residual $\rightarrow$ maturidade.}
\end{principio}

Nenhuma ameaça relevante é considerada controlada apenas porque existe código defensivo. O estado depende de teste reproduzível, evidência válida e risco residual aceito pela autoridade definida.

\section{Contratos e objetos centrais a codificar}

\subsection{\texttt{SystemRegistration}}
Campos mínimos: identificador, nome, domínio, versão, owner, mantenedor alternativo, capacidades, endpoints, contratos suportados, diagnóstico, revisão e datas. Devem ser independentes:
\begin{itemize}
  \item \texttt{registry\_status}: \texttt{registered}, \texttt{active}, \texttt{suspended}, \texttt{retired};
  \item \texttt{maturity\_profile}: perfil avaliado pelo SGE, como \texttt{experimental}, \texttt{shadow} ou estágio posterior.
\end{itemize}
Invariantes: identidade única, capacidade versionada, owner presente e transições auditáveis.

\subsection{\texttt{IntegrationAgreement}}
Campos mínimos: participantes, capacidades, finalidade, escopo, restrições, schemas, revisão, digest, aceites, estado, vigência, política de falha, segurança e evidências. Estados mínimos: \texttt{proposed}, \texttt{validated}, \texttt{accepted}, \texttt{active}, \texttt{suspended}, \texttt{revoked}, \texttt{superseded}.

\subsection{\texttt{IntegrationRun}}
Campos mínimos:
\begin{itemize}
  \item identificador e chave de idempotência;
  \item acordo, revisão, participantes, ator, finalidade e capacidade;
  \item entradas, digests, schemas e transformações;
  \item instantes, tentativas, bloqueios, erros, saídas e evidências;
  \item relações \texttt{retry\_of\_run\_id}, \texttt{reprocesses\_run\_id} e \texttt{supersedes\_run\_id}, quando aplicáveis.
\end{itemize}
O estado de execução deve ser separado da validade posterior:
\begin{itemize}
  \item \texttt{execution\_status}: \texttt{requested}, \texttt{authorized}, \texttt{validating}, \texttt{running};
  \item estados terminais: \texttt{succeeded}, \texttt{failed}, \texttt{blocked}, \texttt{cancelled};
  \item \texttt{validity\_status}: \texttt{valid}, \texttt{invalidated}, \texttt{superseded}.
\end{itemize}
Retry ou reprocessamento cria nova execução vinculada à anterior.

\subsection{\texttt{ProvenanceRecord}}
Deve representar relações como \texttt{wasGeneratedBy}, \texttt{used}, \texttt{wasDerivedFrom}, \texttt{wasAttributedTo} e \texttt{wasInvalidatedBy}. Uma implementação relacional versionada é suficiente no primeiro ciclo se permitir travessia de impacto.

\subsection{\texttt{IntegratedInformation}}
Campos mínimos: identificador, entidades canônicas, tempo, espaço, fontes, correlações, transformação, qualidade, inconsistências, pendências, validade, responsável e proveniência.

\subsection{\texttt{KnowledgeArtifact}}
Campos mínimos: tipo, afirmação, informações sustentadoras, método, evidências, autor ou responsável, confiança, limitações, validade temporal, finalidade, revisões e substituições. Dimensões independentes:
\begin{itemize}
  \item \texttt{artifact\_type}: \texttt{scientific}, \texttt{operational}, \texttt{administrative}, \texttt{methodological}, \texttt{governance};
  \item \texttt{review\_status}: \texttt{draft}, \texttt{under\_review}, \texttt{accepted}, \texttt{rejected};
  \item \texttt{publication\_status}: \texttt{unpublished}, \texttt{published}, \texttt{withdrawn};
  \item \texttt{validity\_status}: \texttt{valid}, \texttt{invalidated}, \texttt{superseded}, \texttt{retracted}.
\end{itemize}
O tipo determina método mínimo e autoridade revisora. Revisão científica é obrigatória apenas para artefatos \texttt{scientific}.

\subsection{\texttt{DecisionUse}}
Relaciona conhecimento, decisão, ator, finalidade, contexto, restrições, data e resultado observado. O registro torna o uso auditável e não atribui decisão autônoma ao SisTer.

\subsection{\texttt{ReflexivityProfile}}
Contrato operacional versionado que declara: escopo; profundidade D0--D5; autoridade A0--A5; modo; tipos de referência; evidências obrigatórias; avaliador e versão; tolerâncias; mapeamento de severidade; efeito de gate; ações corretivas permitidas; owner de autorização; requisito de rollback; retenção; e política de falha. Deve possuir schema, digest, vigência e compatibilidade explícita. Uma execução relevante sem perfil válido produz \texttt{inconclusive} ou bloqueio definido por política, nunca avaliação improvisada.

\subsection{\texttt{ReferenceSnapshot}}
Congela as referências utilizadas em uma avaliação: objetivo e função, processo, acordo e revisão, políticas, critérios de aceitação, perfil de risco, dependências, contexto e instante de vigência. Deve possuir digest e não pode ser reescrito após a avaliação; mudança posterior cria nova referência.

\subsection{\texttt{OperationalAssessment}}
Campos mínimos: identificador, execução/processo avaliado, \texttt{ReflexivityProfile}, \texttt{ReferenceSnapshot}, evidências consideradas, avaliador e versão, instante, resultado, severidade, divergências, limitações, confiança, recomendação, efeito de gate, resposta operacional proposta, impacto sobre operação/maturidade e responsável. Resultados mínimos: \texttt{confirmed}, \texttt{divergent}, \texttt{inconclusive} e \texttt{not\_applicable}. Resultado, severidade, gate e decisão corretiva são dimensões independentes.

\subsection{\texttt{CorrectiveAction}}
Relaciona avaliação, ação proposta, nível de autonomia permitido, autoridade, aprovação, escopo, precondições, efeitos esperados, reversão, estado, evidências de aplicação e avaliação posterior. Estados mínimos: \texttt{proposed}, \texttt{authorized}, \texttt{applied}, \texttt{verified}, \texttt{rolled\_back}, \texttt{rejected} e \texttt{failed}. Mudança estrutural nunca é autorizada apenas pela existência de uma avaliação divergente.


\subsection{\texttt{ReferenceRevisionProposal}}
Registra descoberta que pode exigir alteração de propósito, contrato, política, critério ou método. Campos mínimos: referência atual, avaliações e evidências motivadoras, mudança proposta, impacto, owner competente, consulta/revisão exigida, decisão, nova versão e relação de substituição. O artefato pode ser criado automaticamente, mas sua aprovação pertence ao nível A5 de governança humana.

\section{Rastreabilidade da engenharia}

\subsection{Cadeia normativa}
\begin{center}
\resizebox{0.92\textwidth}{!}{%
\begin{tikzpicture}[
  node distance=4mm,
  t/.style={draw=sisterblue, rounded corners, fill=sisterlightblue, minimum width=8.8cm, minimum height=8mm, align=center, font=\small\bfseries},
  arr/.style={-{Latex[length=2.1mm]}, thick, draw=sisterdarkgray}
]
\node[t] (o) {Objetivo científico};
\node[t, below=of o] (f) {Função sistêmica};
\node[t, below=of f] (p) {Processo};
\node[t, below=of p] (c) {Contrato};
\node[t, below=of c] (u) {Caso de uso};
\node[t, below=of u] (cl) {Componente / unidade de implementação};
\node[t, below=of cl] (te) {Teste};
\node[t, below=of te] (at) {Atestação};
\node[t, below=of at] (ra) {Avaliação reflexiva e divergências};
\node[t, below=of ra] (m) {Indicador, ação autorizada e decisão de maturidade};
\draw[arr] (o)--(f); \draw[arr] (f)--(p); \draw[arr] (p)--(c); \draw[arr] (c)--(u); \draw[arr] (u)--(cl); \draw[arr] (cl)--(te); \draw[arr] (te)--(at); \draw[arr] (at)--(ra); \draw[arr] (ra)--(m);
\draw[arr, bend right=72] (m.west) to node[left, font=\scriptsize, align=right]{correção e\\reavaliação} (p.west);
\end{tikzpicture}}
\end{center}

\subsection{Convenção de identificadores}
\begin{tabularx}{\textwidth}{P{3.0cm}P{3.3cm}Y}
\toprule
\textbf{Elemento} & \textbf{Padrão} & \textbf{Exemplo} \\
\midrule
Objetivo & \texttt{O-nn} & \texttt{O-01} Produzir convergência governada. \\
Função & \texttt{F-nn} & \texttt{F-06} Preservar proveniência. \\
Processo & \texttt{P-nn} & \texttt{P-03} Recepção de dado. \\
Capacidade & \texttt{CAP-domínio-nn} & \texttt{CAP-EXEC-01} Registrar execução. \\
Contrato & ID + versão & \texttt{integration-run/1.0.0}. \\
Caso de uso & \texttt{UC-domínio-nn} & \texttt{UC-EXEC-01} Iniciar integração. \\
Requisito & \texttt{REQ-domínio-nn} & \texttt{REQ-PROV-04} Relacionar derivado à origem. \\
Ameaça & \texttt{TH-domínio-nn} & \texttt{TH-HTTP-01} Parsing e enquadramento inválido. \\
Controle & \texttt{CTRL-domínio-nn} & \texttt{CTRL-HTTP-03} Conter exceções por job. \\
Teste & \texttt{TST-domínio-nn} & \texttt{TST-EXEC-07} Rejeitar acordo suspenso. \\
Evidência & schema + id da execução & \texttt{sister.integration-attestation/1.0.0}. \\
Perfil reflexivo & \texttt{RFP-domínio-nn} & \texttt{RFP-NC-01} Nexo--Compras D2--D3/A1 em shadow. \\
Avaliação & \texttt{ASM-domínio-nn} & \texttt{ASM-REF-01} Comparar execução ao acordo ativo. \\
Ação corretiva & \texttt{ACT-domínio-nn} & \texttt{ACT-REF-02} Suspender integração divergente. \\
Indicador & \texttt{IND-nn} & \texttt{IND-04} Proporção com proveniência válida. \\
\bottomrule
\end{tabularx}

\subsection{Matriz inicial de rastreabilidade}
\begin{landscape}
\scriptsize
\begin{longtable}{P{1.3cm}P{1.3cm}P{1.3cm}P{3.4cm}P{3.5cm}P{3.8cm}P{3.6cm}P{3.0cm}}
\toprule
\textbf{Obj.} & \textbf{Função} & \textbf{Proc.} & \textbf{Capacidade} & \textbf{Caso de uso} & \textbf{Implementação-alvo} & \textbf{Teste/evidência} & \textbf{Indicador} \\
\midrule
\endhead
O-01 & F-01 & P-01 & CAP-FED-01 registrar sistema & UC-FED-01 registrar manifesto & \texttt{FederationService}, \texttt{SystemRepository}, API/CLI & manifesto válido/inválido, duplicidade, atestação de registro & IND-01 sistemas com perfil válido \\
O-01 & F-02 & P-02 & CAP-AGR-01 ativar acordo & UC-AGR-03 ativar após aceite & \texttt{AgreementService}, repositório, política de transição & rejeitar aceite incompleto; recibo e digest & IND-02 integrações com acordo ativo \\
O-01 & F-03 & P-01--P-07 & CAP-IAM-01 autorizar capacidade & UC-IAM-02 avaliar ator/finalidade & \texttt{AuthorizationPolicy}, identidade humana e de sistema & matriz de autorização e decisão registrada & IND-03 operações autorizadas/negadas \\
O-02 & F-04/F-05 & P-03 & CAP-EXEC-01 iniciar execução & UC-EXEC-01 receber dado & \texttt{IntegrationRunService}, validador, adaptador & contrato inativo, schema inválido, idempotência & IND-04 execuções reproduzíveis \\
O-02 & F-06 & P-03--P-07 & CAP-PROV-01 registrar derivação & UC-PROV-01 ligar entrada e saída & \texttt{ProvenanceService}, \texttt{ProvenanceRepository} & travessia origem--derivado e impacto & IND-05 derivados com proveniência válida \\
O-02 & F-07 & P-04 & CAP-INF-01 contextualizar & UC-INF-01 correlacionar necessidade/projeto & \texttt{ContextualizationService}, adaptador Nexo--Compras & relação confirmada, pendência e conflito & IND-06 dados contextualizados \\
O-03 & F-08 & P-05 & CAP-KNW-01 registrar conhecimento & UC-KNW-01 submeter para revisão & \texttt{KnowledgeService}, revisão e repositório & evidência insuficiente, revisão, invalidação & IND-07 artefatos revisados \\
O-03 & F-09 & P-05/P-06 & CAP-EXP-01 publicar & UC-EXP-01 expor conforme escopo & API, CLI, interface e política de exposição & acesso permitido/negado e versão correta & IND-08 usos governados \\
O-04 & F-10 & todos & CAP-OPS-01 diagnosticar & UC-OPS-01 executar doctor federado & contrato de diagnóstico, agregador SGE, CLI & PASS/WARN/BLOCK, sanitização, schema & IND-09 checks resolvidos \\
O-05 & F-11 & P-08 & CAP-REF-01 aplicar perfil e avaliar execução & UC-REF-01 comparar execução e referências & \texttt{ReflexivityService}, \texttt{AssessmentEngine}, repositórios de perfil, referência, avaliação e correção & perfil válido, confirmação/divergência/inconclusão, gate separado, explicação, autorização e verificação & IND-10 avaliações reproduzíveis por perfil e correções eficazes \\
\bottomrule
\end{longtable}
\end{landscape}

\subsection{Rastreabilidade proporcional}
A rastreabilidade obrigatória incide sobre unidades que alterem comportamento observável, contrato, segurança, persistência, integração ou maturidade. Funções auxiliares são relacionadas pelo módulo, requisito e conjunto de testes.

\subsection{Regra de revisão de código}
Nenhuma mudança governada deve ser aprovada sem responder:
\begin{enumerate}
  \item Qual objetivo e função ela atende?
  \item Qual processo e caso de uso são alterados?
  \item Qual contrato, schema ou política limita o comportamento?
  \item Qual teste demonstra o critério de aceitação?
  \item Qual evidência será produzida?
  \item Quais ameaças e fronteiras são afetadas e quais controles se aplicam?
  \item Qual teste negativo e qual evidência demonstram o controle?
  \item Como a execução será comparada às referências e que divergências podem ser produzidas?
  \item Qual ação corretiva é permitida, por quem, com que reversão e verificação posterior?
  \item Qual risco residual, compatibilidade e impacto sobre maturidade existem?
\end{enumerate}

\section{Roadmap de materialização por capacidades}

O roadmap é organizado em trilhas coordenadas. A base segura precede alterações estruturais; a fatia Nexo--Compras orienta a generalização; operação e SGE acompanham os incrementos.

\begin{landscape}
\scriptsize
\begin{longtable}{P{1.4cm}P{2.4cm}P{4.2cm}P{6.0cm}P{5.4cm}P{2.5cm}}
\toprule
\textbf{ID} & \textbf{Trilha} & \textbf{Capacidade} & \textbf{Implementação principal} & \textbf{Gate} & \textbf{Dependência} \\
\midrule
\endhead
I-00 & Base segura & Caracterização e decisão HTTP & Cobrir comportamento atual; registrar ADR sobre manter, endurecer ou substituir o servidor HTTP próprio. & Rotas, códigos, autenticação e comandos preservados; ADR aceita. & -- \\
I-01 & Base segura & Fronteiras e portas mínimas & Extrair transporte, roteamento e composição; definir portas de persistência, identidade, relógio e integrações. & \texttt{main.cpp} reduzido sem mudança de contrato externo. & I-00 \\
I-02 & Fatia de referência & Registro e acordo mínimos & Persistir sistema, capacidade e acordo Nexo--Compras com estados independentes e aceite. & Execução rejeitada sem sistema ativo e acordo válido. & I-01 \\
I-03 & Fatia de referência & \texttt{IntegrationRun} mínima & Criar execução, idempotência, entradas, estado, saída, erro e relações entre execuções. & Uma execução é identificável e reconstruível. & I-02 \\
I-04 & Fatia de referência & Proveniência e informação integrada & Ligar original, execução e transformação; produzir correlação ou pendência Nexo--Compras. & Resultado remonta às entradas e explicita qualidade e pendências. & I-03 \\
I-05 & Fatia de referência & Conhecimento tipificado & Produzir artefato \texttt{operational} ou \texttt{governance}; habilitar revisão conforme tipo, invalidação e substituição. & Artefato revisável e vinculado à informação integrada. & I-04 \\
I-06 & Operação & Cliente \texttt{sisterctl} & Implementar comandos operacionais, \texttt{--json}, modo não interativo e códigos estáveis; separar manutenção. & Fluxo de referência operável sem acesso direto à persistência. & I-02--I-05 \\
I-07 & Operação & Saúde, prontidão e auditoria & Separar \textit{liveness}, \textit{readiness} e degradação; implementar \texttt{doctor}, logs correlacionados e exportação segura. & Dependências e falhas são detectáveis e auditáveis. & I-01; evolui com I-03--I-06 \\
I-08 & SGE & Avaliação em \texttt{shadow} & Criar schemas e checks para acordo, execução, proveniência, informação, conhecimento e CLI desde a primeira fatia. & Cada incremento produz atestação válida sem bloquear inicialmente. & inicia em I-02 \\
I-09 & SGE/operação & Gates vinculantes e recuperação & Ativar bloqueios por risco; implementar retenção, backup, rollback, regressão e resposta a falhas. & Promoção e publicação dependem de evidência; recuperação comprovada. & I-05--I-08 \\
I-10 & Reflexividade & Perfis e avaliação em \texttt{shadow} & Implementar \texttt{ReflexivityProfile}, \texttt{ReferenceSnapshot}, coleta RAII, \texttt{AssessmentEngine} e \texttt{OperationalAssessment}; configurar Nexo--Compras em D2--D3/A1. & A mesma execução e referências reproduzem o mesmo resultado; CLI explica resultado, severidade e gate sem alterar o fluxo. & I-04--I-09 \\
I-11 & Reflexividade & Recomendação, gates e correção limitada & Evoluir perfis aceitos para A2/A3 e, somente em ações locais reversíveis, A4; separar avaliador e executor; persistir autorização, aplicação, rollback e nova avaliação. & Nenhuma ação excede o perfil; bloqueios e correções possuem política, testes de falha, recibo, rollback e verificação de efetividade. & I-10, I-06--I-09 \\
\bottomrule
\end{longtable}
\end{landscape}


\subsection{Trilha de segurança}
\begin{longtable}{P{1.8cm}P{3.2cm}P{5.5cm}P{2.9cm}}
\toprule
\textbf{ID} & \textbf{Entrega} & \textbf{Gate} & \textbf{Dependência} \\
\midrule
\endhead
SEC-GOV-00 & MAES-SisTer e perfil do \texttt{sisterd} & Ativos, fronteiras, classes de ameaça, controles, owners e risco residual estão versionados. & I-00 \\
SEC-00 & Quarentena do transporte & Produção aceita apenas transporte interno; proxies e bootstrap legados falham fechados. & SEC-GOV-00 \\
SEC-01 & Autorização e bootstrap & Políticas negam por padrão; capacidades são explícitas; bootstrap administrativo é local e de uso único. & SEC-00 \\
SEC-01C & Robustez do parser e dos workers & Entrada malformada produz erro controlado; worker e processo sobrevivem; teste de regressão existe. & SEC-01 \\
SEC-01D & Contenção de abuso & Rate limiting não depende da porta, possui memória limitada e combina IP, identidade e limite global. & SEC-01 \\
SEC-02 & Identidade interna assinada & Cookie humano não atravessa fronteira; audiência, validade, capacidade, rotação e falha fechada são testadas. & SEC-01C, SEC-01D \\
SEC-03 & Gateway especializado & TLS, normalização, limites externos, WebSocket, origem confiável e observabilidade de borda estão operacionais. & SEC-02 \\
SEC-SGE-01 & Verificação contínua & Threat IDs, controles, testes e risco residual integram atestações e regressão de maturidade. & inicia em SEC-GOV-00 \\
\bottomrule
\end{longtable}

\subsection{Coordenação das trilhas}
\begin{lstlisting}
I-00 -> I-01 -> I-02 -> I-03 -> I-04 -> I-05
                    |       |       |       |
                    +------ I-08 acompanha em shadow
I-02 -----------------------------> I-06
I-01 -> I-07, ampliado a cada fatia
I-05 + I-06 + I-07 + I-08 -> I-09
I-04 + I-08 + I-09 -> I-10

SEC-GOV-00 -> SEC-00 -> SEC-01 -> SEC-01C + SEC-01D
SEC-01C + SEC-01D -> SEC-02 -> SEC-03
SEC-SGE-01 acompanha toda a trilha
\end{lstlisting}

\section{Requisitos transversais de engenharia}

\subsection{Segurança}
Aplicam-se as obrigações da Seção~\ref{sec:engenharia-seguranca}. Em particular:
\begin{itemize}
  \item nova fronteira, protocolo, parser, proxy, credencial ou exposição exige modelo de ameaças e perfil de segurança;
  \item credenciais e segredos nunca entram em evidências, logs ou repositório;
  \item identidade de sistema é distinta da identidade humana e restrita por audiência e capacidade;
  \item autorização nega por padrão e evolui para capacidade, finalidade e escopo;
  \item transporte possui limites de tamanho, tempo, concorrência, memória, destino e taxa;
  \item falha provocada por entrada não confiável deve ser contida sem encerrar o plano de controle;
  \item evidências são classificadas, sanitizadas, retidas e protegidas contra alteração indevida;
  \item operação destrutiva exige autorização explícita, confirmação adequada, recuperação e auditoria;
  \item avaliação reflexiva e ação corretiva devem respeitar separação de autoridade, menor privilégio, reversibilidade e falha fechada.
\end{itemize}

\subsection{Confiabilidade e resiliência}
\begin{itemize}
  \item idempotência em comandos repetíveis e chaves em execuções;
  \item timeout e política de retry por capacidade, sem repetição cega de operação não idempotente;
  \item degradação explícita: fallback nunca pode parecer dado atual e autorizado;
  \item reprocessamento preserva execução anterior e produz nova relação de proveniência;
  \item regressão automática ou assistida quando evidência, dependência ou contrato perde validade;
  \item ação corretiva com escopo, reversão e verificação posterior, sem apagar a execução divergente.
\end{itemize}

\subsection{Testabilidade}
Cada caso de uso deve possuir testes de unidade das invariantes, testes de contrato, integração de persistência/adaptadores, cenários de falha e pelo menos um teste ponta a ponta por processo. Testes devem produzir resultados consumíveis pelo SGE.

\subsection{Observabilidade}
Logs devem possuir identificador de requisição e, quando aplicável, \texttt{run\_id}, ator, capacidade, contrato, resultado e latência, sem segredos. Métricas mínimas: taxa de erro, tempo por estado, filas, retries, indisponibilidade, pendências, invalidações, reprocessamentos, avaliações por resultado, divergências por severidade e efetividade das correções.

\subsection{Reflexividade operacional}
\begin{itemize}
  \item toda função classificada como relevante deve declarar propósito, \texttt{ReflexivityProfile}, referências avaliáveis, evidências mínimas e owner do critério;
  \item profundidade D0--D5, autoridade A0--A5, efeito de gate e ações permitidas devem ser explícitos e versionados;
  \item toda avaliação deve usar referências congeladas e avaliador versionado, permitindo reprodução e contestação;
  \item ausência ou insuficiência de perfil, referência ou evidência produz \texttt{inconclusive} ou bloqueio definido por política, nunca confirmação implícita;
  \item resultado da avaliação, severidade, efeito de gate, impacto sobre maturidade e decisão corretiva são dimensões independentes;
  \item divergências devem ser explicadas em forma humana e estruturada, com evidências favoráveis e contrárias;
  \item recomendações não são autorizações; ações automáticas limitam-se às previamente aprovadas por política, proporcionais ao risco, auditáveis e reversíveis;
  \item toda correção aplicada deve produzir evidência e ser submetida a nova avaliação de efetividade;
  \item revisão de propósito, contrato, política, critério ou método cria \texttt{ReferenceRevisionProposal}, exige autoridade A5 e não altera retrospectivamente avaliações anteriores;
  \item recursos experimentais do C++26 não podem ser requisito da linha de base até ADR, portabilidade e testes demonstrarem adoção segura.
\end{itemize}

\subsection{Compatibilidade}
Schemas e contratos são versionados. Mudança incompatível exige nova versão, migração ou adaptador explícito. Compatibilidade deve ser comprovada por contrato e teste.

\section{Gates contra \textit{vibe coding}}

\subsection{Definição de pronto para iniciar}
Uma capacidade somente entra em implementação governada quando possui:
\begin{itemize}
  \item problema e objetivo relacionados;
  \item função, processo e responsável;
  \item escopo e fora de escopo;
  \item contrato ou hipótese de contrato;
  \item critérios de aceitação verificáveis;
  \item ameaças aplicáveis, fronteiras afetadas, controles e risco residual esperado;
  \item estratégia mínima de teste negativo, contenção e evidência;
  \item referência de avaliação, resultado esperado e limite de autonomia para eventual correção.
\end{itemize}

\subsection{Definição global de concluído}
Uma capacidade somente é concluída quando:
\begin{itemize}
  \item código e contratos estão versionados;
  \item critérios são satisfeitos por testes reproduzíveis;
  \item evidência estruturada foi produzida e validada;
  \item documentação e rastreabilidade foram atualizadas;
  \item compatibilidade, dados, ameaças, controles e risco residual foram revisados;
  \item testes negativos, contenção, operação, falha, rollback/reprocessamento e observabilidade foram tratados;
  \item owner científico/funcional validou o significado do resultado;
  \item SGE avaliou o perfil aplicável;
  \item quando aplicável, P-08 produziu avaliação reflexiva válida e a correção foi verificada ou formalmente dispensada.
\end{itemize}

\subsection{Sinais de bloqueio}
Devem bloquear integração ou promoção:
\begin{itemize}
  \item implementação sem critério de aceitação;
  \item contrato ausente ou estado não verificável;
  \item \texttt{SKIP} sem motivo estruturado;
  \item fallback indistinguível de resultado real;
  \item alteração direta de banco por interface operacional;
  \item evidência sem revisão, ambiente ou artefato relacionado;
  \item decisão crítica presente apenas em conversa;
  \item nova superfície ou protocolo sem ameaça, controle, teste negativo e owner;
  \item ameaça conhecida e aplicável sem tratamento, exceção expirada ou risco residual aceito;
  \item parser ou fronteira exposta sem contenção de falhas e limites de recursos demonstrados;
  \item mudança de comportamento, contrato, segurança, persistência ou integração sem requisito, capacidade e teste associados;
  \item função relevante sem referência avaliável, evidência suficiente ou owner do critério;
  \item divergência bloqueante ou inconclusão crítica sem decisão, prazo e responsável;
  \item ação corretiva automática sem autorização prévia, reversibilidade ou verificação posterior.
\end{itemize}

\section{Processo de referência: Nexo--Compras}

\subsection{Objetivo demonstrador}
Comprovar que o SisTer transforma uma necessidade de compra em informação contextualizada por projeto e atividade de pesquisa e em artefato de conhecimento operacional ou de governança, sem transferir ao núcleo a autoridade do Nexo ou do Compras.

\subsection{Entradas e autoridades}
\begin{tabularx}{\textwidth}{P{4.0cm}P{4.2cm}Y}
\toprule
\textbf{Oferta} & \textbf{Autoridade} & \textbf{Uso no processo} \\
\midrule
Necessidade de compra & Compras & Origem operacional, prioridade, estado, itens e justificativa. \\
Projeto e atividade científica & Nexo & Finalidade científica, desafio, responsável e vínculo institucional. \\
Acordo e contratos técnicos & Relação/SisTer & Autoriza capacidades, finalidade, schemas, restrições e evidências. \\
Execução e proveniência & SisTer & Registra como, quando e sob qual revisão a contextualização ocorreu. \\
Validação do artefato & Autoridade definida pelo tipo & Define significado, confiança, limites e uso admissível; revisão científica aplica-se apenas ao tipo \texttt{scientific}. \\
\bottomrule
\end{tabularx}

\subsection{Resultado mínimo do ciclo}
\begin{enumerate}
  \item acordo Nexo--Compras ativo e versionado;
  \item \texttt{IntegrationRun} criada por comando/API;
  \item dados originais preservados por referência/digest;
  \item validações e transformações registradas;
  \item necessidade vinculada a projeto/atividade ou marcada como pendente com razão;
  \item \texttt{IntegratedInformation} produzida;
  \item ao menos um \texttt{KnowledgeArtifact} do tipo \texttt{operational} ou \texttt{governance} submetido à revisão;
  \item invalidação de uma entrada identifica os derivados afetados;
  \item SGE publica atestação do processo;
  \item o perfil \texttt{RFP-NC-01} aplica profundidade D2--D3, autoridade A1 e efeito \texttt{shadow};
  \item P-08 compara execução e resultados ao propósito, ao acordo e aos critérios aplicáveis, publicando \texttt{OperationalAssessment};
  \item uma divergência controlada produz recomendação, decisão autorizada, correção reversível e nova avaliação;
  \item todo o ciclo pode ser executado e consultado pelo \texttt{sisterctl}.
\end{enumerate}

\subsection{Critérios de aceitação ponta a ponta}
\begin{longtable}{P{1.5cm}P{9.0cm}P{4.1cm}}
\toprule
\textbf{ID} & \textbf{Critério} & \textbf{Evidência} \\
\midrule
\endhead
AC-01 & Execução sem acordo ativo é rejeitada antes de consultar dados de domínio. & Decisão de autorização e status \texttt{blocked}. \\
AC-02 & A mesma chave idempotente não duplica efeitos. & Duas solicitações, um efeito e referência à execução original. \\
AC-03 & Cada entrada possui produtor, schema, digest e referência ao original. & Registro de entrada e validação. \\
AC-04 & Correlação inexistente não é inventada; permanece pendente com causa explícita. & \texttt{IntegratedInformation} com pendência. \\
AC-05 & Correlação confirmada referencia identificadores dos dois domínios. & Relação e proveniência reproduzíveis. \\
AC-06 & Alteração do projeto ou contrato identifica informações e conhecimentos afetados. & Consulta de impacto. \\
AC-07 & Artefato de conhecimento registra tipo, método, evidências, confiança, limitações, responsável e autoridade revisora. & Schema válido e revisão compatível com o tipo. \\
AC-08 & CLI apresenta resultado humano e JSON equivalente, com códigos estáveis. & Teste de contrato do CLI. \\
AC-09 & Evidência não contém credenciais, tokens ou URLs com segredos. & Check de sanitização. \\
AC-10 & SGE distingue \texttt{PASS}, \texttt{SHADOW}, \texttt{WARN} e \texttt{BLOCK}. & Atestação validada por schema. \\
AC-11 & A execução é avaliada sob perfil D2--D3/A1 contra propósito, acordo, políticas e critérios congelados, produzindo confirmação, divergência ou inconclusão explicável sem alterar o fluxo. & \texttt{ReflexivityProfile}, \texttt{ReferenceSnapshot} e \texttt{OperationalAssessment} válidos. \\
AC-12 & Resultado da avaliação, severidade, efeito de gate e resposta proposta permanecem campos independentes e coerentes com o perfil. & Teste de schema e cenários \texttt{divergent+warn}, \texttt{divergent+block} e \texttt{inconclusive}. \\
AC-13 & Divergência de teste não altera contrato, método ou código automaticamente; ação autorizada é reversível, auditada e reavaliada. & \texttt{CorrectiveAction}, evidência de aplicação/rollback, avaliação posterior e eventual \texttt{ReferenceRevisionProposal}. \\
\bottomrule
\end{longtable}

\section{Responsabilidades e governança do trabalho}

\begin{tabularx}{\textwidth}{P{4.1cm}Y}
\toprule
\textbf{Papel} & \textbf{Responsabilidade final} \\
\midrule
Coordenação científica & Finalidade científica, critérios de conhecimento científico, método, limites, referências de avaliação e validação do significado quando o tipo exigir. \\
Arquitetura e governança & Funções, processos, fronteiras, ADRs, contratos estruturais e rastreabilidade. \\
Core / \texttt{sisterd} & Casos de uso transversais, estado, persistência, autorização, proveniência e APIs. \\
CLI / \texttt{sisterctl} & Experiência operacional, automação, saída estruturada, códigos e segurança de comando. \\
Owners dos subsistemas & Capacidades, contratos de domínio, diagnóstico, testes e autoridade dos dados. \\
Integração e qualidade & Adaptadores, testes de contrato, ponta a ponta, falhas e evidências. \\
SGE-SisTer & Schemas de atestação e avaliação reflexiva, checks, comparação, explicação de divergências, gates, maturidade, regressão e publicação. \\
Segurança/operação & MAES-SisTer, classes de ameaça, controles, risco residual, segredos, observabilidade, níveis de autonomia corretiva, continuidade, recuperação, advisories e resposta a incidentes. \\
Owner funcional/científico da referência & Propósito, critério, tolerâncias, severidade e condições que distinguem confirmação, divergência e inconclusão. \\
\bottomrule
\end{tabularx}

Uma pessoa pode acumular papéis, mas cada capacidade e cartão deve possuir um único responsável final. O owner técnico não deve ser a única pessoa capaz de operar, avaliar ou recuperar a capacidade.


\subsection{Comunicação e adoção pela equipe}
A versão 1.4 deve ser comunicada antes do início da implementação reflexiva porque altera vocabulário, fronteiras de autoridade, critérios de pronto e desenho de interfaces. A comunicação mínima deve:
\begin{enumerate}
  \item apresentar a distinção entre reflexividade do SisTer e reflexão do C++;
  \item explicar as escalas D0--D5 e A0--A5 e a separação entre resultado, gate e ação;
  \item declarar que o piloto Nexo--Compras começa em D2--D3/A1/\texttt{shadow};
  \item informar que nenhuma alteração de contrato, método, política ou código será automática;
  \item abrir a ADR-REF-01 para estabilizar schema do \texttt{ReflexivityProfile}, interfaces C++ e critérios de evolução a A2--A4;
  \item separar recursos C++20/23 da linha de base e experimentos C++26 do laboratório.
\end{enumerate}

A comunicação não representa ordem para implementar imediatamente todo o modelo. Ela alinha a linguagem da equipe, impede soluções incompatíveis e define o experimento mínimo que produzirá evidência para a próxima decisão.

\section{Manutenção desta especificação}

\subsection{Quando atualizar}
Atualização é obrigatória quando houver:
\begin{itemize}
  \item nova função sistêmica ou processo finalístico;
  \item mudança de fronteira ou autoridade;
  \item novo objeto central ou transição de estado;
  \item mudança incompatível em contrato estruturante;
  \item alteração na cadeia de rastreabilidade;
  \item nova superfície, protocolo ou classe de ameaça que altere garantia estrutural;
  \item perda de controle, evidência ou risco residual incompatível com a maturidade;
  \item constatação de que o código real divergiu do modelo-alvo;
  \item alteração de propósito, referência, critério, tolerância, severidade ou nível de autonomia corretiva.
\end{itemize}

\subsection{Como lidar com divergências}
Divergência entre esta especificação e o código deve ser classificada:
\begin{itemize}
  \item \textbf{defeito}: código viola requisito vigente;
  \item \textbf{dívida}: implementação parcial aceita temporariamente;
  \item \textbf{descoberta}: experimento revelou necessidade de revisar a especificação;
  \item \textbf{exceção}: desvio temporário, com risco, owner e expiração;
  \item \textbf{obsolescência}: requisito perdeu finalidade e deve ser removido por decisão registrada.
  \item \textbf{divergência operacional}: a execução não correspondeu à referência vigente e exige classificação de severidade, decisão e eventual correção.
  \item \textbf{inconclusão}: faltam referência, evidência ou capacidade avaliativa; deve gerar pendência explícita e não pode ser tratada como conformidade.
\end{itemize}

\subsection{Critério de autoridade}
Divergência entre especificação e código exige decisão versionada, atualização dos testes e nova evidência. Nenhuma divergência pode ser resolvida implicitamente.

\section{Síntese executiva}

O snapshot verificado apresenta servidor, autenticação, gestão de usuários, consultas PostgreSQL, proxies, clientes de integração e exposição de maturidade. Ainda não materializa execução governada, proveniência causal, informação integrada, conhecimento tipificado ou avaliação reflexiva da própria execução.

\begin{principio}[Síntese adotada]
\centering
\textbf{Objetivo $\rightarrow$ função $\rightarrow$ processo $\rightarrow$ capacidade $\rightarrow$ contrato $\rightarrow$ caso de uso $\rightarrow$ implementação $\rightarrow$ teste $\rightarrow$ evidência $\rightarrow$ maturidade.}
\end{principio}

A evolução deve preservar o comportamento atual, estabelecer fronteiras e implementar fatias verticais verificáveis. O processo Nexo--Compras será a primeira comprovação: acordo válido, \texttt{IntegrationRun}, original preservado, proveniência, informação contextualizada, artefato operacional ou de governança, invalidação, avaliação pelo SGE e avaliação reflexiva regida por \texttt{ReflexivityProfile} D2--D3/A1 em \texttt{shadow}. Somente evidência acumulada e decisão arquitetural podem elevar autoridade ou efeito operacional.

\begin{principio}[Síntese de segurança]
\centering
\textbf{O SisTer limita a superfície que decide manter, herda o conhecimento acumulado sobre ameaças aplicáveis e somente promove capacidades cujos controles foram demonstrados por evidências.}
\end{principio}

\begin{principio}[Síntese de reflexividade]
\centering
\textbf{O SisTer deve explicar o que fez, avaliar se permaneceu coerente com aquilo que deveria fazer e tornar governável o que precisa acontecer em seguida.}
\end{principio}

\appendix
\section{Backlog mínimo rastreável}
\begin{longtable}{P{1.5cm}P{3.6cm}P{5.7cm}P{2.7cm}}
\toprule
\textbf{ID} & \textbf{Entrega} & \textbf{Pronto quando} & \textbf{Dependência} \\
\midrule
\endhead
SEC-GOV-00 & MAES-SisTer/1.0 & Ativos, fronteiras, ameaças, controles, testes, risco residual, owners e processo de atualização estão versionados. & -- \\
SEC-01C & Robustez do transporte & Entrada HTTP inválida não encerra worker ou processo; corpus negativo e teste de sobrevivência aprovados. & SEC-GOV-00 \\
SEC-01D & Contenção de abuso & Rate limiting usa IP sem porta, identidade e limite global; memória e expiração são limitadas. & SEC-GOV-00 \\
SEC-02 & Identidade interna & Asserção assinada, curta e restrita substitui cookie e headers não verificáveis entre sistemas. & SEC-01C, SEC-01D \\
SEC-03 & Gateway especializado & Borda assume TLS, normalização, quotas, WebSocket, origem confiável e proteção contra requisição lenta. & SEC-02 \\
SEC-SGE-01 & Checks de segurança & Ameaças, controles, testes e risco residual produzem atestação e podem regredir maturidade. & SEC-GOV-00; incremental \\
ENG-00 & Testes de caracterização & Rotas, códigos, autenticação e comandos atuais estão cobertos. & -- \\
ADR-HTTP-01 & Decisão do servidor HTTP & Manutenção, endurecimento ou substituição possui critérios, riscos, owner e testes. & ENG-00 \\
ENG-01 & Extração incremental & Transporte e roteamento são separados; portas mínimas são definidas; contratos externos permanecem estáveis. & ENG-00, ADR-HTTP-01 \\
FED-01 & Registro persistente & Estado operacional e perfil de maturidade são independentes; API e CLI registram, consultam e retiram. & ENG-01 \\
AGR-01 & Ciclo de acordos & Propor, aceitar, ativar, suspender e revogar funcionam com invariantes. & FED-01 \\
EXEC-01 & \texttt{IntegrationRun} v1 & Estado de execução, validade, entradas, resultados, bloqueios, idempotência e vínculos entre execuções são persistidos. & AGR-01 \\
PROV-01 & Proveniência v1 & Original, execução, transformação e saída estão ligados; impacto é consultável. & EXEC-01 \\
INF-01 & Informação Nexo--Compras & Contextualização produz confirmação ou pendência explícita. & PROV-01 \\
KNW-01 & Conhecimento tipificado & Artefato operacional ou de governança possui revisão, publicação e validade independentes. & INF-01 \\
CLI-01 & Cliente operacional & Fluxo de referência funciona por CLI humana e JSON, sem acesso direto à persistência. & FED-01--KNW-01 \\
OPS-01 & Saúde e \texttt{doctor} & \textit{Liveness}, \textit{readiness}, degradação e dependências são consultáveis. & ENG-01; incremental \\
SGE-01 & Avaliação do processo & Checks acompanham cada fatia em \texttt{shadow} e produzem atestações válidas. & inicia em AGR-01 \\
SGE-02 & Gates vinculantes & Promoção, regressão e publicação dependem das evidências aplicáveis. & CLI-01, OPS-01, SGE-01 \\
ADR-REF-01 & Níveis e perfis de reflexividade & Schema, profundidades, autoridades, gates, ações, interfaces C++ e critérios de promoção do piloto estão aceitos. & EFE-SisTer/1.4 \\
REF-00 & \texttt{ReflexivityProfile} v1 & Perfil versionado declara D2--D3/A1/\texttt{shadow}, referências, evidências, avaliador, tolerâncias, gates, owner, retenção e falha. & ADR-REF-01 \\
REF-01 & Avaliação reflexiva v1 & Referências são congeladas; execução Nexo--Compras produz confirmação, divergência ou inconclusão reproduzível e explicável, sem alterar o fluxo. & REF-00, EXEC-01, PROV-01, INF-01, SGE-01 \\
REF-02 & Ação corretiva governada & Recomendação, autorização, aplicação, rollback e verificação posterior são persistidos; nenhuma mudança estrutural ocorre autonomamente. & REF-01, CLI-01, OPS-01, SGE-02 \\
REF-LAB-01 & Laboratório C++26 & Contratos, reflexão estática e \texttt{std::execution} são prototipados atrás de interfaces estáveis, sem dependência da linha de base. & ADR-REF-01, ENG-01 \\
\bottomrule
\end{longtable}


\section{Referências normativas e históricas de segurança}
\label{sec:referencias-seguranca}
\begin{itemize}
  \item IETF. \textit{RFC 9112 -- HTTP/1.1}. Regras de sintaxe, enquadramento de mensagens, tratamento de erro, request smuggling e response splitting. \url{https://www.rfc-editor.org/info/rfc9112/}
  \item NIST. \textit{SP 800-218 -- Secure Software Development Framework (SSDF) Version 1.1}. Práticas de preparação, proteção, produção segura e resposta a vulnerabilidades. \url{https://csrc.nist.gov/pubs/sp/800/218/final}
  \item Apache HTTP Server Project. \textit{Apache HTTP Server 2.4 vulnerabilities}. Histórico de vulnerabilidades corrigidas e classes recorrentes de falha. \url{https://httpd.apache.org/security/vulnerabilities_24.html}
  \item Apache HTTP Server Project. \textit{Security Tips}. Limites de requisição, timeouts, privilégios e mitigação de negação de serviço. \url{https://httpd.apache.org/docs/2.4/en/misc/security_tips.html}
  \item MITRE. \textit{Common Weakness Enumeration (CWE)}. Vocabulário para classificação de fraquezas de software. \url{https://cwe.mitre.org/}
\end{itemize}


\section{Referências de engenharia C++ para a reflexividade}
\label{sec:referencias-cpp}
Estas referências sustentam a direção técnica, mas não substituem a matriz de compiladores, ADRs e testes do projeto:
\begin{itemize}
  \item WG21. \textit{P0323R12 -- std::expected}. \url{https://www.open-std.org/jtc1/sc22/wg21/docs/papers/2022/p0323r12.html}
  \item WG21. \textit{P1208R0 -- source\_location para C++20}. \url{https://www.open-std.org/JTC1/SC22/WG21/docs/papers/2018/p1208r0.html}
  \item WG21. \textit{P2175R0 -- cancelamento componível em operações assíncronas}. \url{https://open-std.org/jtc1/sc22/wg21/docs/papers/2020/p2175r0.html}
  \item WG21. \textit{P2996R13 -- Reflection for C++26}. \url{https://www9.open-std.org/JTC1/SC22/WG21/docs/papers/2025/p2996r13.html}
  \item WG21. \textit{N5015 -- Working Draft, Standard for Programming Language C++}. \url{https://www.open-std.org/jtc1/sc22/wg21/docs/papers/2025/n5015.html}
  \item WG21. \textit{Índices de propostas 2024--2026}, incluindo evolução de \texttt{std::execution} e contratos. \url{https://www.open-std.org/jtc1/sc22/wg21/docs/papers/2026/}
\end{itemize}

\section{Checklist de uma mudança governada}
\begin{enumerate}
  \item Há um identificador e um responsável?
  \item O objetivo, a função e o processo estão declarados?
  \item As fronteiras de autoridade permanecem corretas?
  \item O contrato, o tipo de artefato e a compatibilidade estão definidos?
  \item Há critérios de aceitação e casos de falha?
  \item Quais ameaças conhecidas ou plausíveis se aplicam à fronteira alterada?
  \item Os controles preventivos, detectivos, de contenção e recuperação possuem owner?
  \item A unidade de implementação está na camada correta e respeita as dependências?
  \item Os testes são reproduzíveis e incluem entrada hostil, falha e sobrevivência quando aplicável?
  \item A evidência é estruturada, sanitizada e associada à revisão?
  \item Qual referência congelada permitirá avaliar a execução e distinguir confirmação, divergência e inconclusão?
  \item Que ação corretiva é permitida, quem a autoriza, como é revertida e como sua efetividade será verificada?
  \item O risco residual e o impacto sobre maturidade, operação e segurança foram avaliados?
  \item A decisão relevante saiu da conversa e entrou no repositório?
\end{enumerate}

\end{document}
